% concatenated from main.tex, macros.tex, sections/01_introduction.tex, sections/02_related_work.tex, sections/03_setting.tex, sections/04_main_result.tex, figures/client_averaging_schematic.tex, sections/05_mechanism.tex, sections/06_proof_overview.tex, sections/07_numerical_study.tex, sections/08_limitations_conclusion.tex, appendix/A_scope_and_inputs.tex, appendix/B_exact_recursions.tex, appendix/C_control_and_coarse_moments.tex, appendix/D_local_second_moment.tex, appendix/E_mean_and_global_moments.tex, appendix/F_sharp_second_moment.tex, appendix/G_local_higher_moments_taylor.tex, appendix/H_main_proof.tex, appendix/I_uniformity_and_semantics.tex, appendix/J_numerical_details.tex
% Single-root arXiv build: manuscript and appendices are compiled in one
% LaTeX document so all cross-references are native.
\documentclass[11pt]{article}
\def\TSPREVIEW{1}
\def\ARXIVCOMBINED{1}

\usepackage[T1]{fontenc}
\usepackage[utf8]{inputenc}
\usepackage{lmodern}
\usepackage{microtype}
\usepackage{amsfonts,amssymb,mathtools,bm,mathrsfs}
\usepackage{amsthm}
\usepackage{aliascnt}
\usepackage{geometry}
\usepackage{iftex,ifpdf}
\usepackage[font=footnotesize]{subcaption}
\usepackage[colorlinks=true,allcolors=blue!45!black]{hyperref}
\usepackage[nameinlink,capitalise,noabbrev]{cleveref}
\usepackage{booktabs,array,longtable,tabularx}
\usepackage[numbers,sort&compress]{natbib}
\usepackage{graphicx}
\usepackage{float}
\usepackage{placeins}
\usepackage{xcolor}
\definecolor{clientblue}{HTML}{2F6B9A}
\definecolor{controlorange}{HTML}{B45F06}
\definecolor{biasred}{HTML}{8E3B46}
\definecolor{inkgray}{HTML}{4F5965}
\definecolor{softgray}{HTML}{F2F4F7}
\usepackage{enumitem}
\usepackage{tikz}
\usetikzlibrary{arrows.meta,positioning,fit,calc}
\usepackage{fancyvrb}
\usepackage{url}

\ifpdf
  \DeclareGraphicsExtensions{.pdf,.png,.jpg,.eps}
\else
  \DeclareGraphicsExtensions{.eps}
\fi

\setlist[itemize]{topsep=3pt,itemsep=2pt,parsep=1pt}
\setlist[enumerate]{topsep=3pt,itemsep=2pt,parsep=1pt}
\renewcommand{\arraystretch}{1.12}
\providecommand{\creflastconjunction}{, and~}
\newcommand{\email}[1]{\href{mailto:#1}{#1}}
\allowdisplaybreaks

\newtheorem{assumption}{Assumption}
\newtheorem{theorem}{Theorem}[section]

\newaliascnt{proposition}{theorem}
\newtheorem{proposition}[proposition]{Proposition}
\aliascntresetthe{proposition}

\newaliascnt{lemma}{theorem}
\newtheorem{lemma}[lemma]{Lemma}
\aliascntresetthe{lemma}

\newaliascnt{corollary}{theorem}
\newtheorem{corollary}[corollary]{Corollary}
\aliascntresetthe{corollary}

\theoremstyle{definition}

\newaliascnt{definition}{theorem}
\newtheorem{definition}[definition]{Definition}
\aliascntresetthe{definition}

\newaliascnt{remark}{theorem}
\newtheorem{remark}[remark]{Remark}
\aliascntresetthe{remark}

\newaliascnt{proofnote}{theorem}
\newtheorem{proofnote}[proofnote]{Proof note}
\aliascntresetthe{proofnote}

\AddToHook{env/theorem/begin}{\crefalias{section}{theorem}}
\AddToHook{env/proposition/begin}{\crefalias{section}{proposition}}
\AddToHook{env/lemma/begin}{\crefalias{section}{lemma}}
\AddToHook{env/corollary/begin}{\crefalias{section}{corollary}}
\AddToHook{env/definition/begin}{\crefalias{section}{definition}}
\AddToHook{env/remark/begin}{\crefalias{section}{remark}}
\AddToHook{env/proofnote/begin}{\crefalias{section}{proofnote}}

\crefname{assumption}{Assumption}{Assumptions}
\Crefname{assumption}{Assumption}{Assumptions}
\crefname{theorem}{Theorem}{Theorems}
\Crefname{theorem}{Theorem}{Theorems}
\crefname{proposition}{Proposition}{Propositions}
\Crefname{proposition}{Proposition}{Propositions}
\crefname{lemma}{Lemma}{Lemmas}
\Crefname{lemma}{Lemma}{Lemmas}
\crefname{corollary}{Corollary}{Corollaries}
\Crefname{corollary}{Corollary}{Corollaries}
\crefname{remark}{Remark}{Remarks}
\Crefname{remark}{Remark}{Remarks}
\crefname{proofnote}{Proof note}{Proof notes}
\Crefname{proofnote}{Proof note}{Proof notes}

\newcommand{\E}{\mathbb{E}}
\newcommand{\R}{\mathbb{R}}
\newcommand{\Var}{\operatorname{Var}}
\newcommand{\Cov}{\operatorname{Cov}}
\newcommand{\Scaf}{\textnormal{\textsc{SCAFFOLD}}}
\newcommand{\Oh}{O_H}
\newcommand{\smallscale}{s_{\gamma,N}}
\newcommand{\remscale}{r_{\gamma,N}}
\newcommand{\fourthscale}{v_{\gamma,N}}
\newcommand{\bstat}{b_{\gamma,N,H}}
\newcommand{\QH}{Q_{\gamma,N,H}}
\newcommand{\Ubar}{\bar U}
\newcommand{\eps}{\varepsilon}
\newcommand{\barq}{\overline{q}}
\newcommand{\barf}{\overline{f'}}
\newcommand{\bareps}{\bar\varepsilon}
\newcommand{\delt}{\delta}
\newcommand{\law}{\mathcal{L}}
\newcommand{\statezero}{\mathcal{X}_0}

\hypersetup{
  colorlinks=true,
  linkcolor=blue!50!black,
  citecolor=blue!50!black,
  urlcolor=blue!50!black
}

\geometry{a4paper,margin=1in}

\title{Beyond Client Averaging: A Client-Independent Second-Order Stationary-Bias Component in Stochastic SCAFFOLD}
\author{Yi-Ping Tang and Guan-Ju Peng
\thanks{Both authors are with the Graduate Institute of Data Science and Information Computing,
National Chung Hsing University, Taichung, Taiwan.
Email (Yi-Ping Tang): \email{s4109053021@gmail.com};
Email (Guan-Ju Peng): \email{gjpeng@email.nchu.edu.tw};
ORCID (Guan-Ju Peng): \url{https://orcid.org/0000-0001-5508-9485}.}
}
\date{}

\ifpdf
\hypersetup{
  hypertexnames=false,
  pdftitle={Beyond Client Averaging: A Client-Independent Second-Order Stationary-Bias Component in Stochastic SCAFFOLD},
  pdfauthor={Yi-Ping Tang; Guan-Ju Peng},
  pdfsubject={Stationary bias in stochastic SCAFFOLD},
  pdfkeywords={federated learning, SCAFFOLD, stationary bias, stochastic approximation, control variates}
}
\fi

\begin{document}
\maketitle
\begin{abstract}

Existing constant-step analysis of stochastic \Scaf{} identifies a leading $O(\gamma/N)$ stationary mean bias and shows that higher-order bias can persist as the client count increases, but does not identify the first client-independent contribution at coefficient level. For full-participation stochastic \Scaf{} with one-dimensional homogeneous clients, fixed local-step count $H$, and bounded additive gradient noise, we prove, uniformly over $N\ge2$,

$$
\begin{aligned}
\mathbb{E}_{\pi_{\gamma,N,H}}[x]-x^\star
={}&
-\frac{f'''(x^\star)\sigma^2}{4f''(x^\star)^2}\frac{\gamma}{N}\\
&-
\frac{f'''(x^\star)\sigma^2}{12f''(x^\star)}
\frac{(H-1)(5H-1)}{H}\gamma^2
+O_H\!\left(\frac{\gamma^2}{N}+\gamma^3\right).
\end{aligned}
$$

Hence client averaging suppresses the leading $O(\gamma/N)$ bias but does not remove the client-independent $O(\gamma^2)$ component when its coefficient is nonzero.

The mechanism is indirect: although the direct control contribution cancels pathwise in the linear global average, the controls still alter within-round local trajectories and their second moments. Fresh gradient noise and persistent control fluctuations therefore generate local second-moment corrections that nonquadratic curvature converts into stationary mean bias. The coefficient vanishes for quadratic objectives. Numerical experiments are consistent with the predicted coefficient, its persistence as client count increases, and the stated joint remainder. The result is restricted to the one-dimensional homogeneous fixed-$H$ setting.

\end{abstract}

\section{Introduction}
\label{sec:intro}

Federated learning trades communication for local computation: each participating client performs several stochastic-gradient steps before the server averages the resulting updates. This design can substantially reduce synchronization, but the local trajectories may drift toward client-specific directions. \Scaf{} addresses this problem by attaching control variates to local updates, enabling clients to track a common optimization direction more closely \citep{karimireddy2020scaffold}. Most analyses study finite-time convergence. Under a constant step size, however, stochastic iterates continue to fluctuate around the optimum, and their long-run behavior is described by an invariant distribution whose mean need not equal the minimizer.

Recent constant-step analysis of stochastic \Scaf{} makes the usual ``more clients help'' intuition precise and identifies its higher-order limit. \citet{mangold2025scaffold} establish stationarity, a client-number linear speed-up up to higher-order terms, and a higher-order stationary-bias component that is not eliminated by increasing the client count. In the scalar homogeneous specialization, their leading stationary mean-bias layer is of order $\gamma/N$, where $\gamma$ is the step size and $N$ is the number of participating clients. Still, their analysis leaves the client-independent higher-order contribution unresolved at the coefficient level. We ask whether client averaging suppresses that first client-independent contribution and, if not, which term first carries it and what mechanism generates it.

Our answer is negative in this setting. For full-participation stochastic \Scaf{} with one-dimensional homogeneous clients, fixed local-step count $H$, sufficiently small constant step size $\gamma$, and bounded additive gradient noise, we prove
\begin{equation}
\label{eq:intro-main}
\E_{\pi_{\gamma,N,H}}[x]-x^\star
= -\frac{\tau\sigma^2}{4a^2}\frac{\gamma}{N}
-\frac{\tau\sigma^2}{12a}\frac{(H-1)(5H-1)}{H}\gamma^2
+O_H\!\left(\frac{\gamma^2}{N}+\gamma^3\right),
\end{equation}
where $a=f''(x^\star)$ and $\tau=f'''(x^\star)$. The first displayed layer decreases as the client count increases; the second does not. The uniform remainder further shows that, after the known $\gamma/N$ layer is removed, the normalized residual converges to the displayed $\gamma^2$ coefficient along every joint sequence $N\to\infty$ and $\gamma\to0$.

\paragraph{Why this matters.}
Client averaging reduces both stationary fluctuations and the known leading mean-bias layer, which suggests that increasing participation should continually improve long-run accuracy. Equation~\eqref{eq:intro-main} identifies a limit of that mechanism within the expansion: when its coefficient is nonzero, the client-independent $\gamma^2$ component is not reduced by further averaging after the $\gamma/N$ layer becomes smaller. Increasing $N$ alone cannot reduce this component; decreasing the step size does so, and coefficient-aware bias correction is a natural direction for future work.

Let $Y$ denote a local iterate obtained by sampling uniformly over clients $c$ and local-step indices $h=0,\ldots, H-1$ under stationarity. The exact stationary mean-balance identity gives $\E[f'(Y)]=0$, and a local Taylor expansion then gives the heuristic
\[
0=\E[f'(Y)]
\approx a\,\E[Y]+\frac{\tau}{2}\E[Y^2],
\qquad
\E[Y]\approx-\frac{\tau}{2a}\E[Y^2].
\]
Thus, the second-moment-to-mean intuition follows directly from exact stationary balance plus a local Taylor expansion: nonquadratic curvature converts a second-moment correction into a mean shift. In \Scaf{}, the controls cancel out in the linear client average, but they still affect the intermediate local trajectories. Fresh gradient noise and persistent control fluctuations therefore modify the second moments of local trajectories before curvature converts them into stationary bias; \cref{fig:mechanism} summarizes this indirect channel.

\begin{figure}[H]
\centering
\includegraphics[width=\linewidth]{figures/scaffold_bias_mechanism.png}
\caption{Mechanism behind the client-independent stationary-bias layer. Linear control cancellation at the server does not erase the controls' effect on intermediate local trajectories. Fresh noise and persistent control fluctuations alter local second moments, which nonquadratic curvature converts into a mean shift.}
\label{fig:mechanism}
\end{figure}

The contributions are as follows.
\begin{itemize}
  \item \textbf{A limit to client averaging at second order.} We identify an explicit client-independent $\gamma^2$ stationary-bias coefficient, with a remainder uniform jointly in $(\gamma,1/N)$ for fixed $H$.
  \item \textbf{A trajectory-level mechanism inside \Scaf{}.} We trace the coefficient to two local second-moment sources: fresh stochastic-gradient noise and persistent control fluctuations.
  \item \textbf{Targeted numerical checks.} We test persistence with increasing client count, convergence toward the predicted coefficient, and behavior when stochasticity or nonquadratic curvature is removed.
\end{itemize}

The scalar analysis is informative beyond one dimension because it isolates three ingredients that also occur in higher-dimensional dynamics: persistent control second moments, within-round local trajectories, and nonlinear curvature. Their interaction would involve covariance--third-derivative contractions; establishing such an extension remains open.

\paragraph{Claim boundaries.}
The theorem concerns full participation, one-dimensional homogeneous clients, bounded additive noise, and fixed $H$. It is uniform in $N$ and small $\gamma$, but not as $H\to\infty$, and it does not identify the complete finite-$N$ coefficient of $\gamma^2$. The source decomposition is internal to \Scaf{} and is not a theorem-level comparison with FedAvg. The crossover of the two displayed layers is an asymptotic scale comparison, not a monotonicity statement for the total finite-step bias. The numerical experiments are finite-setting consistency checks. Multidimensional, heterogeneous-client, state-dependent-noise, and debiasing extensions remain open.

\Cref{sec:related} positions the result. \Cref{sec:setting} connects the standard \Scaf{} controls to the zero-sum parametrization used in the proof. \Cref{sec:main-result,sec:mechanism} state the theorem and derive its mechanism, while \cref{sec:proof-overview,sec:numerics} summarize the proof and numerical evidence.

\section{Related work}
\label{sec:related}

\paragraph{Federated averaging and local SGD.}
Federated Averaging (FedAvg) introduced iterative local model averaging as a communication-efficient primitive for federated learning \citep{mcmahan2017fedavg}; the broader optimization and systems landscape is surveyed by \citet{kairouz2021federated}. Its optimization core is closely related to LocalSGD, in which workers take several stochastic gradient steps before synchronization. Early theory established that LocalSGD can retain the statistical rate of minibatch SGD while communicating less \citep{stich2019localsgd}, and subsequent analyses clarified the distinct roles of identical and heterogeneous client objectives \citep{khaled2020localsgd} and the regimes in which local updates do or do not improve on minibatch SGD \citep{woodworth2020localsgd}.

\paragraph{SCAFFOLD and control correction.}
\Scaf{} augments local updates with control variates to correct client drift \citep{karimireddy2020scaffold}. More recent work sharpens the finite-time picture: \citet{luo2025revisiting} obtain improved LocalSGD and \Scaf{} rates under gradient/Hessian similarity and higher-order smoothness conditions, while \citet{mangold2025sharper} give a refined quadratic analysis for arbitrary local-step counts. These results explain when local computation and control correction improve finite-time optimization, but they do not characterize the mean of the invariant distribution generated by constant-step stochastic \Scaf{}.

\paragraph{Constant-step stochastic optimization and stationary laws.}
Under suitable stability and regularity conditions, constant-step stochastic-gradient methods are naturally studied in terms of an invariant distribution. This viewpoint appears in diffusion approximations of constant-step SGD \citep{mandt2017sgd} and in rigorous small-step characterizations of stationary stochastic-gradient dynamics \citep{dieuleveut2020bridging,chen2022stationary}. In particular, \citet{chen2022stationary} obtain Gaussian/Lyapunov characterizations of scaled stationary laws for smooth strongly convex SGD and related stochastic-approximation models under their conditions, while \citet{dieuleveut2020bridging} derive stationary/asymptotic moment expansions and use Richardson--Romberg extrapolation to reduce step-size bias. The same perspective has recently been developed for federated and decentralized algorithms. \citet{mangold2025fedavg} analyze constant-step FedAvg through its Markov structure, characterize stationary bias and variance, derive a first-order bias expansion, and construct a federated Richardson--Romberg correction. \citet{mangold2025scaffold} extend this framework to stochastic \Scaf{}, proving existence and geometric convergence of a stationary law, client-number variance reduction, and a first-order stationary mean-bias expansion whose scalar homogeneous leading layer is $O(\gamma/N)$. \citet{versini2026dsgd} obtain related first-order stationary bias and variance expansions for decentralized SGD, separating stochastic, heterogeneity, and network effects.

\paragraph{Nonlinearity and higher-order stationary bias.}
The mechanism studied here belongs to a broader class of noise--nonlinearity interactions in constant-step stochastic approximation. \citet{allmeier2024bias} allow Markovian noise to depend on the iterate, prove an $O(\alpha)$ stationary-bias bound, and refine the Polyak--Ruppert time-averaged bias to $\alpha V+O(\alpha^2)$. \citet{huo2024collusion} show that Markovian memory and nonlinearity can interact to create distinct components of the constant-step stationary bias. These results establish that persistent stochastic fluctuations can be converted into mean shifts by nonlinear dynamics. They do not, however, determine the algorithm-specific higher-order coefficient generated by \Scaf{}'s local control recursion.

\paragraph{Position of the present result.}
The closest starting point is \citet{mangold2025scaffold}. We use their stationary-law existence and geometric convergence, coarse iterate-moment bounds, leading control second moment, and known $O(\gamma/N)$ bias coefficient. Mangold et al. establish that stochastic \Scaf{} retains a higher-order stationary bias that is not eliminated by increasing the client count. Still, their analysis does not identify the client-independent contribution at the coefficient level. In the one-dimensional homogeneous fixed-$H$ setting, we resolve this structure explicitly: the first client-independent term is $N^0\gamma^2$, with a closed-form coefficient that decomposes into fresh within-round local-gradient noise and persistent \Scaf{} control second moments. The result is complementary to the quadratic theory of \citet{mangold2025sharper}: quadratic objectives support sharper arbitrary-$H$ convergence analysis, but $f'''\equiv0$ removes the second-moment-to-mean conversion isolated here. The source decomposition is internal to \Scaf{}; it is not a theorem-level comparison with FedAvg, which would require a matched second-order FedAvg expansion.

\section{Setting and exact identities}
\label{sec:setting}

We deliberately use a homogeneous scalar regime, removing client heterogeneity as a confounding source of bias and isolating the stochastic effect of the \Scaf{} control recursion itself. All $N$ clients participate in every communication round. Even with identical objectives, the stochastic client controls remain nontrivial because they are continually refreshed from noisy local trajectories. The unique minimizer is translated to $x^\star=0$, all $N\ge2$ clients share the same objective $f_c=f$, and the number of local steps $H\ge2$ is fixed. The objective satisfies $f\in C^5(\R)$ and
\[
0<\mu\le f''(y)\le L<\infty,
\]
with bounded third through fifth derivatives. We write
\[
a:=f''(0)>0,\qquad \tau:=f'''(0),\qquad \kappa:=f''''(0).
\]
At local step $h$ of client $c$, the stochastic gradient is
\[
\nabla F_{c,h}(y)=f'(y)+\eps_{c,h},
\]
where the fresh noises are independent across clients, local steps, and communication rounds, independent of the round-start state, satisfy $\E[\eps_{c,h}]=0$ and $\E[\eps_{c,h}^2]=\sigma^2$, and are almost surely bounded. No symmetry assumption is imposed.

At the start of a communication round, the server model is $x$, and every client initializes at $\theta_{c,0}=x$. Each client then takes $H$ corrected stochastic-gradient steps. We first connect the parametrization used below to the standard \Scaf{} controls.

\begin{center}
\fbox{\begin{minipage}{0.93\linewidth}
\textbf{Equivalent full-participation SCAFFOLD round.}
Let $c_{\rm srv}$ be the server control and $c_c$ the control stored by client $c$, with
$c_{\rm srv}=N^{-1}\sum_c c_c$. Standard \Scaf{} uses the correction
$c_{\rm srv}-c_c$ in each local step. Define
\[
\xi_c:=c_{\rm srv}-c_c,
\qquad \frac1N\sum_{c=1}^N\xi_c=0.
\]
With full participation, global mixing parameter one, and the standard Option-II control refresh \citep{karimireddy2020scaffold},
\[
c_c^+=c_c-c_{\rm srv}+\frac{x-\theta_{c,H}}{\gamma H},
\qquad
c_{\rm srv}^+=c_{\rm srv}+\frac1N\sum_{j=1}^N(c_j^+-c_j).
\]
Using $x^+=N^{-1}\sum_j\theta_{j,H}$ gives
\[
\xi_c^+=c_{\rm srv}^+-c_c^+
=\xi_c+\frac{\theta_{c,H}-x^+}{\gamma H}.
\]
Thus the zero-sum variables $\{\xi_c\}$ are an exact reparametrization of the usual server/client controls in the regime analyzed here.
\end{minipage}}
\end{center}

Using this parametrization, the local, server, and control updates are
\begin{align}
\theta_{c,h+1}
&=\theta_{c,h}-\gamma\{f'(\theta_{c,h})+\xi_c+\eps_{c,h+1}\},
\qquad \theta_{c,0}=x,\label{eq:local-update-main}\\
x^+&=\frac1N\sum_{c=1}^N\theta_{c,H},\label{eq:global-update-main}\\
\xi_c^+&=\xi_c+\frac{\theta_{c,H}-x^+}{\gamma H}.\label{eq:control-update-main}
\end{align}
Summing Equation~\eqref{eq:control-update-main} over clients shows directly that the zero-sum condition is preserved.

Two exact identities expose the distinction between global cancellation and local influence. First, averaging the unrolled local recursion gives
\begin{equation}
\label{eq:global-cancel-main}
x^+-x
=-\gamma\sum_{h=0}^{H-1}\overline{f'(\theta_h)}
-\gamma\sum_{h=0}^{H-1}\bar\eps_{h+1},
\end{equation}
where bars denote client averages. The controls cancel pathwise from this linear average. Second, stationarity gives
\begin{equation}
\label{eq:stationary-balance-main}
0=\sum_{h=0}^{H-1}\frac1N\sum_{c=1}^N\E[f'(\theta_{c,h})].
\end{equation}
These identities contain the central tension of the paper. Equation~\eqref{eq:global-cancel-main} rules out a direct linear control contribution to the global update, but it does not erase the effect of $\xi_c$ on the intermediate trajectories $\theta_{c,h}$. Equation~\eqref{eq:stationary-balance-main} is where those trajectory-level changes re-enter the stationary mean through the nonlinear map $f'$. Formal assumptions and the earlier stationary results used below are stated in Appendix~\ref{sec:scope}.

\section{Main result: a client-independent second-order layer}
\label{sec:main-result}

Let $\pi_{\gamma, N, H}$ denote the unique stationary law of the global/control process, whose existence for sufficiently small $\gamma$ follows from \cref{prop:anchor-inputs}, and define
\[
\bstat:=\E_{\pi_{\gamma,N,H}}[x]-x^\star.
\]

\begin{theorem}[Uniform joint stationary-bias expansion]
\label{thm:main}
Under the assumptions of \cref{sec:setting}, including full participation, for fixed $H\ge2$ there exist $\gamma_{0,H}>0$ and $C_H<\infty$, independent of $N$ and $\gamma$, such that for every $N\ge2$ and $0<\gamma\le\gamma_{0,H}$,
\begin{equation}
\label{eq:main-expansion}
\boxed{
\bstat
=-\frac{\tau\sigma^2}{4a^2}\frac{\gamma}{N}
-\frac{\tau\sigma^2}{12a}\frac{(H-1)(5H-1)}{H}\gamma^2
+R_{\gamma, N , H}
}
\end{equation}
with
\begin{equation}
\label{eq:main-remainder}
\boxed{
|R_{\gamma,N,H}|\le C_H\left(\frac{\gamma^2}{N}+\gamma^3\right).
}
\end{equation}
\end{theorem}

The client-independent $\gamma^2$ coefficient is
\begin{equation}
\label{eq:B20-main}
\boxed{
B_{20}^{\mathrm{SCAF}}(H)
=-\frac{f'''(x^\star)\sigma^2}{12f''(x^\star)}
\frac{(H-1)(5H-1)}{H}.
}
\end{equation}
The subscript $20$ records the bivariate order $\gamma^2N^0$: second order in the step size and independent of the client count.

\subsection{Interpretation for client averaging and local computation}
\label{sec:main-interpretation}

\paragraph{More clients suppress only the leading displayed layer.}
The known $O(\gamma/N)$ contribution decreases at rate $1/N$, whereas the displayed $\gamma^2$ contribution is independent of $N$. Thus, within the theorem's stationary and small-step regime, client averaging suppresses the leading bias layer but does not remove the client-independent second-order layer when its coefficient is nonzero. When $\tau\sigma^2\neq0$, equating the magnitudes of the two displayed terms gives
\[
N_\times
=\frac{3H}{a(H-1)(5H-1)}\,\gamma^{-1}.
\]
The two displayed coefficients then have the same sign. This is a crossover scale for the two displayed asymptotic layers, not an exact finite-step phase transition or a monotonicity statement for the full bias; \cref{fig:client-averaging-schematic} visualizes this scale comparison.

\begin{figure}[H]
\centering
\begin{tikzpicture}[
  x=1.36cm,
  y=0.93cm,
  axis/.style={-{Latex[length=2.3mm]},line width=0.75pt,draw=inkgray},
  guide/.style={densely dotted,line width=0.75pt,draw=inkgray!70},
  decreasing/.style={line width=1.55pt,draw=clientblue},
  constant/.style={line width=1.55pt,draw=controlorange,dashed},
  label/.style={font=\small,align=center},
  note/.style={font=\footnotesize,align=center,text=inkgray},
  crossover/.style={font=\footnotesize,align=left,text=inkgray,
                    fill=white,rounded corners=1pt,inner sep=2.5pt}
]
  % Axes: the vertical scale is intentionally schematic.
  \draw[axis] (0.55,0.25) -- (8.45,0.25);
  \draw[axis] (0.55,0.25) -- (0.55,4.48);
  \node[font=\small,below=8pt] at (4.50,0.25)
    {number of participating clients $N$};
  \node[font=\small,rotate=90,anchor=south] at (0.12,2.40)
    {schematic magnitude of displayed bias layers};

  % The two displayed asymptotic components are shown separately.
  \draw[decreasing]
    plot[smooth,domain=0.72:8.05,samples=100]
    (\x,{3.20/(\x+0.35)+0.20});
  \draw[constant] (0.72,1.10) -- (8.05,1.10);

  % Displayed-layer crossover; this is not a total-bias transition curve.
  \draw[guide] (3.20,0.25) -- (3.20,3.82);
  \fill[inkgray] (3.20,1.10) circle (1.15pt);
  \node[note,fill=white,inner sep=1pt] at (3.20,0.53) {$N_{\times}$};
  \node[crossover,anchor=west] at (3.55,3.62)
    {\textbf{Displayed-layer crossover}\\[-1pt]
     $\displaystyle N_{\times}
       =\frac{3H}{a(H-1)(5H-1)}\,\gamma^{-1}$};

  % Component labels.
  \node[label,fill=white,rounded corners=1pt,inner sep=2pt,
        text=clientblue!85!black] at (1.85,2.55)
    {$O(\gamma/N)$ component};
  \node[note,fill=white,inner sep=1pt] at (1.85,2.15)
    {decreases with $N$};

  \node[label,fill=white,rounded corners=1pt,inner sep=2pt,
        text=controlorange!90!black] at (6.05,1.58)
    {client-independent $O(\gamma^2)$ component\\[-1pt]
     \footnotesize does not decay with $N$};
\end{tikzpicture}

\caption{Schematic magnitude of the two displayed terms in \cref{thm:main}. The $O(\gamma/N)$ component decreases with $N$, whereas the client-independent $O(\gamma^2)$ component does not. When both are nonzero, equating their displayed coefficients gives the marked crossover scale. It is not an exact finite-step threshold or a curve for the total bias.}
\label{fig:client-averaging-schematic}
\end{figure}

\paragraph{The coefficient records a local-computation effect.}
For fixed problem parameters, the factor
\[
\frac{(H-1)(5H-1)}{H}
\]
is increasing over integers $H\ge2$, while the theorem treats each $H$ as fixed. The coefficient therefore links the client-independent stationary effect to computation performed inside each communication round.

\paragraph{Quadratic models hide the mechanism.}
When the objective is quadratic, $f'''(x^\star)=0$ and both displayed coefficients vanish. A purely quadratic analysis can therefore characterize stability and convergence sharply while remaining blind to the second-moment-to-mean conversion isolated here.

\subsection{Precise joint-limit meaning}
The uniform remainder gives the exact asymptotic interpretation. After subtracting the known $\gamma/N$ layer and normalizing by $\gamma^2$,
\begin{equation}
\label{eq:normalized-remainder}
\left|
\frac{\bstat+\frac{\tau\sigma^2}{4a^2}\frac{\gamma}{N}}{\gamma^2}
-B_{20}^{\mathrm{SCAF}}(H)
\right|
\le C_H\left(\frac1N+\gamma\right).
\end{equation}
Hence, the normalized residual converges to the value in Equation~\eqref{eq:B20-main} along any joint sequence $N\to\infty$, $\gamma\to0$, without a relative-rate condition. Terms of order $\gamma^2/N$ remain in the remainder, so this identifies the client-independent coefficient but not the complete fixed-$N$ second-order expansion.

\section{Mechanism: how controls survive linear cancellation}
\label{sec:mechanism}

\Cref{fig:mechanism} gives the qualitative picture. The coefficient follows from making the two sources of its fluctuations precise.

\paragraph{Linear aggregation.}
Under the zero-sum control state, Equation~\eqref{eq:global-cancel-main} shows that the client-average control term cancels pathwise from the global update. There is no direct linear control term shifting the server model.

\paragraph{Within-round local dynamics.}
Each $\xi_c$ nevertheless changes the intermediate trajectory
\[
\theta_{c,0},\theta_{c,1},\ldots,\theta_{c,H},
\]
and hence the points at which subsequent stochastic gradients are evaluated. The nonlinear local updates act along these control-dependent paths before the server observes their average.

\paragraph{Two second-moment sources.}
Let
\[
Q_{\gamma,N,H}:=\frac1N\sum_{c=1}^N\E[\xi_c^2]
\]
be the stationary average control second moment. In the homogeneous additive-noise setting,
\begin{equation}
\label{eq:QH-main}
Q_{\gamma,N,H}=\frac{\sigma^2}{H}\left(1-\frac1N\right)+O_H(\gamma).
\end{equation}
The controls therefore retain nonzero fluctuations as the number of clients grows. For the change in the local second moment relative to the round-start server model,
\[
\bar U_h:=\frac1N\sum_{c=1}^N\bigl(\E[\theta_{c,h}^2]-\E[x^2]\bigr),
\]
we obtain
\begin{equation}
\label{eq:U-main}
\bar U_h
=\gamma^2\bigl(h\sigma^2+h^2Q_{\gamma,N,H}\bigr)
+O_H\!\left(\frac{\gamma^2}{N}+\gamma^3\right).
\end{equation}
The first displayed term is generated by fresh within-round gradient noise. The persistent round-start controls generate the second. Because Equation~\eqref{eq:QH-main} approaches $\sigma^2/H$ as $N$ grows, both sources contribute at the client-independent $\gamma^2$ scale.

\paragraph{Nonlinear conversion.}
Expanding the gradient near the optimum,
\[
f'(y)=ay+\frac{\tau}{2}y^2+O(y^3),
\]
and inserting the local moments into the exact stationary balance in Equation~\eqref{eq:stationary-balance-main} converts the second-moment correction into a mean shift. The resulting coefficient separates into
\begin{align}
B_{20}^{\mathrm{local}}(H)
&=-\frac{\tau\sigma^2}{4a}(H-1),\label{eq:Blocal-main}\\
B_{20}^{\mathrm{ctrl}}(H)
&=-\frac{\tau\sigma^2}{12a}\frac{(H-1)(2H-1)}{H},\label{eq:Bctrl-main}\\
B_{20}^{\mathrm{SCAF}}(H)
&=B_{20}^{\mathrm{local}}(H)+B_{20}^{\mathrm{ctrl}}(H).\label{eq:Bsum-main}
\end{align}
The first component records direct local-gradient noise; the second records the additional second moment carried by the \Scaf{} controls. Their sum is the coefficient in Equation~\eqref{eq:B20-main}.

\section{Proof overview}
\label{sec:proof-overview}

The coefficient suggested by a formal Taylor expansion is not automatically identifiable. The stationary balance contains the global second moment, within-round local second-moment corrections, and higher signed moments. Any one of these quantities could, in principle, carry an additional client-independent $\gamma^2$ contribution. The proof must therefore isolate the two intended local sources of fluctuation and rule out all competing channels at the same scale.

The argument proceeds in the order shown in the dependency map in \cref{fig:proof-map} of Appendix~\ref{app:main-proof}. Coarse moment bounds first place the global iterate, controls, and local displacements on scales that are uniform in $N$. The local recursion is then expanded sharply enough to expose the two terms in Equation~\eqref{eq:U-main}: fresh noise and the round-start control second moment. Before these terms can be converted to the stationary mean, the proof controls the global fourth moment and the signed third moment, ensuring that cubic and quartic Taylor contributions remain within the target remainder.

The key identifiability step is the sharp global second moment
\begin{equation}
\label{eq:sharpM2-overview}
\E[x^2]=\frac{\sigma^2}{2a}\frac{\gamma}{N}
+O_H\!\left(\frac{\gamma^2}{N}+\gamma^3\right).
\end{equation}
A coarser $O(\gamma/N+\gamma^2)$ estimate would leave open an additional client-independent $\gamma^2$ term in $\E[x^2]$, which would contaminate the coefficient attributed to local trajectories. Establishing Equation~\eqref{eq:sharpM2-overview} rules out that competing source.

The delicate covariance in this step couples fresh round noise with the nonlinear increment. A generic Cauchy--Schwarz bound is too loose because it loses the extra factor $1/N$ required by the joint remainder. The proof uses a coordinate-replacement argument that sets one fresh-noise coordinate to zero. Only one client's within-round path changes, and the subsequent client average contributes an additional $1/N$ sensitivity. Summing over all noise coordinates then recovers the required scale.

With the global second moment sharpened and the local cubic and quartic terms placed inside the remainder, the exact stationary balance reduces to
\[
\bstat
=-\frac{\tau}{2a}\E[x^2]
-\frac{\tau}{2aH}\sum_{h=0}^{H-1}\bar U_h
+O_H\!\left(\frac{\gamma^2}{N}+\gamma^3\right).
\]
Substituting Equations~\eqref{eq:sharpM2-overview} and~\eqref{eq:U-main}, then summing $h$ and $h^2$, yields \cref{thm:main}. The full proof begins in Appendix~\ref{sec:scope} and is assembled in Appendix~\ref{app:main-proof}.

\section{Numerical checks of the predicted effect}
\label{sec:numerics}

The numerical study follows the paper's main claim in the same order: first, whether the debiased statistic remains nonzero as the client count grows; second, whether it approaches the predicted coefficient as $(\gamma,1/N)\to(0,0)$; and third, whether the observations are consistent with removing the effect when stochasticity or nonquadratic curvature is absent.

We use the smooth, strongly convex objective
\begin{equation}
\label{eq:sim-objective}
f'(x)=x+\frac12\log\cosh(x),
\end{equation}
for which
\[
\frac12\le f''(x)=1+\frac12\tanh(x)\le\frac32,
\qquad a=1,\qquad \tau=\frac12.
\]
The additive noise is independent Rademacher noise with variance one. The predicted client-independent coefficient becomes
\begin{equation}
\label{eq:sim-B}
B_{20}(H)=-\frac{(H-1)(5H-1)}{24H},
\end{equation}
so $B_{20}(2)=-0.1875$ and $B_{20}(4)=-0.59375$.

We estimate the stationary mean both directly and through the exact mean-balance identity. Unless stated otherwise, $\hat b_{\gamma, N, H}$ in the main text denotes the mean-balance estimate; the direct sample mean is retained as a diagnostic. The primary statistic removes the known leading layer and normalizes the residual:
\begin{equation}
\label{eq:Csim}
C_{\gamma, N, H}
:=\frac{\hat b_{\gamma,N,H}+\frac18\frac{\gamma}{N}}{\gamma^2}.
\end{equation}
Thus, the theorem predicts that this normalized residual should approach a nonzero constant as $N$ grows and $\gamma$ decreases.
By \cref{thm:main},
\begin{equation}
\label{eq:Csim-prediction}
C_{\gamma,N,H}=B_{20}(H)+O_H\!\left(\frac1N+\gamma\right).
\end{equation}
Each setting uses eight independent chains, and uncertainty is reported by 95\% Student-$t$ intervals across chain estimates. The complete burn-in, batching, estimator agreement, and stability protocol are given in Appendix~\ref{sec:numerical-supp}.

\subsection{Does the component survive increasing client count?}
\label{sec:num-client}

We fix $H=2$ and $\gamma=1/24$ and increase $N\in\{8,16,32,64,128\}$. After removing the known $\gamma/N$ contribution, the normalized residual approaches a nonzero plateau near the predicted value $-0.1875$; see \cref{fig:nonvanishing}. A regression in $1/N$ gives large-$N$ intercept $-0.189145$ with 95\% bootstrap interval $[-0.190398,-0.188035]$. This is the numerical check that most directly mirrors the title claim: averaging suppresses the leading layer but does not remove the client-independent component.

\begin{figure}[H]
\centering
\includegraphics[width=0.55\linewidth]{figures/client_count_scaling.png}
% \caption{Client-count persistence at $H=2$ and $\gamma=1/24$. The normalized residual approaches a nonzero plateau; larger $N$ lies to the left because the horizontal axis is $1/N$.}
\caption{Client-count persistence at $H=2$ and $\gamma=1/24$. The normalized residual approaches a nonzero plateau; larger $N$ lies to the right on the horizontal axis.}
\label{fig:nonvanishing}
\end{figure}

\subsection{Does the normalized residual approach the coefficient?}
\label{sec:num-coefficient}

Along the joint path $N=1/\gamma$, Equation~\eqref{eq:Csim-prediction} permits an $O_H(\gamma)$ error. The coarse-grid $H=4$ offset is therefore compatible with finite-step effects allowed by the theorem; the smaller-step experiment tests whether the discrepancy contracts at the predicted order. For $H=2$, the initial grid $\gamma\in\{1/12,1/16,1/24,1/32,1/48\}$ gives extrapolated intercept $-0.187214$ with 95\% bootstrap interval $[-0.189250,-0.185175]$, consistent with $B_{20}(2)$.

The same coarse grid at $H=4$ showed a measurable finite-step offset, so we ran a separately specified smaller-step study with fresh chains at
\[
(\gamma,N)\in\{(1/64,64),(1/96,96),(1/128,128),(1/192,192)\}.
\]
The normalized statistic moves toward $B_{20}(4)$, while the scaled error
$\{C_{\gamma,1/\gamma,4}-B_{20}(4)\}/\gamma$ remains between $1.38$ and $1.52$; see \cref{fig:H4-confirm} in Appendix~\ref{sec:numerical-supp}. A linear fit over these four settings yields an intercept of $-0.592953$ with a 95\% bootstrap interval of $[-0.595526,-0.590236]$. The smaller-step figure, initial discrepancy, follow-up criteria, and numerical table are retained in Appendix~\ref{sec:numerical-supp}.

\subsection{Do the required ingredients matter?}
\label{sec:num-removal}

The supplementary negative controls remove both ingredients from the mechanism. With deterministic gradients, the stationary-mean intervals contain zero. With a quadratic objective, where $f'''(x^\star)=0$, direct-mean intervals also contain zero. The settings, intervals, and an additional $H=1$ boundary check are reported in Appendix~\ref{sec:numerical-supp}.

Together, the experiments are consistent with three distinct implications of the expansion: persistence as the client count increases, convergence toward the predicted coefficient, and dependence on stochastic, nonquadratic local dynamics.

\FloatBarrier

\section{Conclusion}
\label{sec:conclusion}
In the homogeneous scalar fixed-$H$ setting, stochastic \Scaf{} has two stationary-bias layers: the known $O(\gamma/N)$ contribution, which client averaging suppresses, and a client-independent $O(\gamma^2)$ contribution that is unaffected by the client count when its coefficient is nonzero. Fresh gradient noise and persistent control fluctuations generate the latter through local second moments, and nonquadratic curvature converts it into a stationary mean shift.

Control cancellation is only linear: the controls vanish from the server average but still change local trajectories. Numerically, the residual forms a client-count plateau, approaches the coefficient along smaller-step paths, and disappears without stochasticity or nonlinearity. In higher dimensions, the scalar second-moment-to-mean conversion would be replaced by contractions between stationary covariance structure and the third-derivative tensor, while client heterogeneity may introduce additional weighted moment channels absent from the present homogeneous argument. Multidimensional analysis, client heterogeneity, and coefficient-based bias cancellation remain open.


\clearpage
\begingroup
\raggedright
\bibliographystyle{plainnat}
\bibliography{references}
\endgroup

\clearpage
\appendix
\section{Scope, notation, and imported results}
\label{sec:scope}

\subsection{Theorem scope}

We translate the optimum to the origin and work only in the following restricted setting.

\begin{assumption}[One-dimensional homogeneous fixed-$H$ setting]
\label{ass:scope}
Let $d=1$, let $N\ge 2$, and fix an integer $H\ge 2$. All $N$ clients participate in every communication round and have the same objective $f_c=f$. The unique minimizer is translated to $x^\star=0$, so $f'(0)=0$. The integer $H$ is fixed: no bound in this document is claimed to be uniform as $H\to\infty$.
\end{assumption}

\begin{assumption}[Strong convexity, smoothness, and higher regularity]
\label{ass:smooth}
There exist $0<\mu\le L<\infty$ such that
\[
  \mu\le f''(y)\le L,\qquad y\in\R.
\]
Moreover, $f\in C^5(\R)$ and, for $j=3,4,5$, there are finite constants $K_j$ such that
\[
  \sup_{y\in\R}|f^{(j)}(y)|\le K_j.
\]
Write
\[
  a:=f''(0)>0,\qquad \tau:=f'''(0),\qquad \kappa:=f''''(0).
\]
\end{assumption}

\begin{assumption}[Bounded additive fresh gradient noise]
\label{ass:noise}
At local step $h$ of client $c$, the stochastic gradient is
\[
  \nabla F_{c,h}(y)=f'(y)+\eps_{c,h}.
\]
The variables $\{\eps_{c,h}\}$ are independent across clients, local steps, and communication rounds, independent of the round-start state, and satisfy
\[
  \E[\eps_{c,h}]=0,\qquad
  \E[\eps_{c,h}^2]=\sigma^2,\qquad
  |\eps_{c,h}|\le B_\eps\quad\text{a.s.}
\]
No symmetry assumption is imposed. In particular, $\E[\eps^3]$ may be nonzero.
\end{assumption}

\begin{remark}[Sample-loss representation of additive noise]
The additive stochastic-gradient oracle in \cref{ass:noise} can equivalently be represented by the sample loss
\[
  F_{c,h}(y)=f(y)+\eps_{c,h}y.
\]
Indeed,
\[
  \nabla F_{c,h}(y)=f'(y)+\eps_{c,h},
  \qquad
  F_{c,h}''(y)=f''(y).
\]
Hence each sample loss inherits the same $L$-smoothness bound as $f$. This verifies the sample-wise smoothness condition required by the imported SCAFFOLD results from \citet{mangold2025scaffold}.
\end{remark}

\begin{assumption}[Zero-sum control state]
\label{ass:zerosum}
The SCAFFOLD state is restricted to the invariant subspace
\[
  \statezero
  :=\left\{(x,\xi_1,\ldots,\xi_N)\in\R^{N+1}:
  \sum_{c=1}^N\xi_c=0\right\}.
\]
Equivalently, the initial controls satisfy $\sum_c\xi_c^0=0$.
\end{assumption}

\begin{remark}[Role of the zero-sum state]
The control sum is conserved exactly by the SCAFFOLD control update; see \cref{lem:exact-recursion}. If the initial sum were nonzero, the global recursion would retain a deterministic linear control drift. The earlier stationary analysis used below is formulated on the corresponding zero-sum state space. This restriction is therefore part of the theorem interface, not merely a proof convenience.
\end{remark}

\subsection{Stationary convention and scale notation}

Throughout the proof, $(x,\Xi)=(x,\xi_1,\ldots,\xi_N)$ denotes a round-start state distributed according to the stationary law $\pi_{\gamma, N, H}$, and all expectations also include the fresh noises drawn during the next communication round. The communication-round index is suppressed.

\begin{table}[H]
\centering
\small
\caption{Core notation used in the main text and appendices.}
\label{tab:core-notation}
\begin{tabularx}{\linewidth}{@{}>{\raggedright\arraybackslash}p{0.22\linewidth}X@{}}
\toprule
Symbol & Meaning\\
\midrule
$x$ & Round-start server model; under stationarity, the global coordinate of $\pi_{\gamma,N,H}$.\\
$\theta_{c,h}$ & Local iterate of client $c$ after $h$ local steps, with $\theta_{c,0}=x$.\\
$\xi_c$ & Zero-sum control correction $c_{\rm srv}-c_c$ used in the local update.\\
$Q_{\gamma,N,H}$ & Stationary average control second moment, $N^{-1}\sum_c\E[\xi_c^2]$.\\
$\bar U_h$ & Local second-moment correction, $N^{-1}\sum_c\{\E[\theta_{c,h}^2]-\E[x^2]\}$.\\
$B_{20}^{\mathrm{SCAF}}(H)$ & Client-independent $\gamma^2N^0$ coefficient in the stationary mean-bias expansion.\\
\bottomrule
\end{tabularx}
\end{table}

Define the three scales
\begin{equation}
\label{eq:scales}
  \smallscale:=\frac{\gamma}{N}+\gamma^2,
  \qquad
  \remscale:=\frac{\gamma^2}{N}+\gamma^3=\gamma\smallscale,
  \qquad
  \fourthscale:=\frac{\gamma^2}{N^2}+\gamma^3.
\end{equation}
The notation $O_H(\cdot)$ means that the hidden constant may depend on fixed $H$ and the fixed problem/noise parameters, but not on $N$ or $\gamma$.

\subsection{Earlier stationary results used in the analysis}

The following inputs are taken from \citet{mangold2025scaffold}; they are not new claims of this document.

\begin{table}[H]
\centering
\small
\caption{Imported-result interface. The middle column records how the assumptions underlying the earlier results are satisfied in the present setting; the last column states how each input is used here.}
\label{tab:imported-interface}
\begin{tabularx}{\linewidth}{@{}>{\raggedright\arraybackslash}p{0.25\linewidth}>{\raggedright\arraybackslash}p{0.34\linewidth}X@{}}
\toprule
Imported result & Assumptions matched here & Role in the present analysis\\
\midrule
Stationarity and geometric $W_2$ convergence
& \Cref{ass:scope,ass:smooth,ass:noise,ass:zerosum}, with the additive oracle represented by the sample loss above
& Defines $\pi_{\gamma,N,H}$ and permits stationary expectation identities.\\
Coarse global/local iterate moments
& The same interface, including bounded noise and sample-wise smoothness
& Supplies the initial moment scales used in the mean and higher-moment bootstraps.\\
Control second-moment expansion
& One-dimensional homogeneous additive-noise specialization of \cref{ass:scope,ass:smooth,ass:noise,ass:zerosum}
& Supplies the persistent control second moment entering the local correction.\\
Leading $O(\gamma/N)$ bias layer
& Scalar homogeneous specialization under \cref{ass:scope,ass:smooth,ass:noise,ass:zerosum}
& Provides the established leading layer that the present uniform expansion refines.\\
\bottomrule
\end{tabularx}
\end{table}

\begin{proposition}[Earlier stationary inputs]
\label{prop:anchor-inputs}
Under the matched assumptions summarized in \cref{tab:imported-interface} and sufficiently small step size, \citet{mangold2025scaffold} establish the following facts.
\begin{enumerate}[label=\textup{(L\arabic*)}]
  \item \textbf{Stationarity.} The joint global/control process is a time-homogeneous Markov chain with a unique stationary distribution and geometric convergence in $W_2$; see \citet[Theorem~4.2]{mangold2025scaffold}.
  \item \textbf{Coarse moments.} The stationary global/local iterates have coarse second moments of order $O_H(\gamma)$, the control second moments are $O_H(1)$, and Appendix~B of \citet{mangold2025scaffold} provides coarse fourth/sixth-moment estimates for the global/local iterates. We use those iterate moment bounds, but we do \emph{not} rely on an imported control sixth-moment bound.
  \item \textbf{Control second moment.} In the one-dimensional homogeneous additive-noise specialization of \citet[Lemma~5.1 and Appendix~E.2]{mangold2025scaffold},
  \begin{equation}
  \label{eq:anchor-QH}
    \QH:=\frac1N\sum_{c=1}^N\E[\xi_c^2]
    =\frac{\sigma^2}{H}\left(1-\frac1N\right)+O_H(\gamma).
  \end{equation}
  In this homogeneous setting, the ideal controls satisfy $\xi_c^\star=-f'(0)=0$, so the squared deviation from the ideal control is the raw second moment displayed above.
  \item \textbf{Known first bias layer.} \citet[Theorem~5.3]{mangold2025scaffold} prove a leading stationary bias of order $\gamma/N$ and leaves a higher-order remainder $O(\gamma^2H+\gamma^{3/2})$. In the present scalar homogeneous specialization, the leading coefficient is
  \begin{equation}
  \label{eq:anchor-bias-leading}
    -\frac{\tau\sigma^2}{4a^2}\frac{\gamma}{N}.
  \end{equation}
\end{enumerate}
\end{proposition}

\begin{remark}[Earlier remainder bound versus the present question]
The terms $O(\gamma^2H)$ and $O(\gamma^{3/2})$ in the earlier theorem are upper remainder bounds. They do not, by themselves, prove the existence of a nonzero $\gamma^2$ coefficient or a nonzero $\gamma^{3/2}$ coefficient. The purpose of the present proof is to resolve the client-number-independent $N^0\gamma^2$ coefficient in the restricted setting above.
\end{remark}

\begin{remark}[Higher-moment envelope versus variance]
The sixth-moment noise envelope in the earlier analysis (often denoted by a symbol such as $\sigma_\star^2$) is not identified with the variance $\sigma^2$ in this document. Whenever an imported sixth-moment bound is used, bounded noise under \cref{ass:noise} supplies the required envelope through $B_\eps$. This distinction is essential for constant bookkeeping.
\end{remark}

\section{Exact SCAFFOLD identities}
\label{sec:exact}

All identities in this section are pathwise and precede any small-step expansion. The one-round update below is the stochastic \Scaf{} recursion specialized to \cref{ass:scope,ass:noise}; it serves as the complete algorithmic definition used in the analysis \citep{mangold2025scaffold}.

\begin{lemma}[Exact local/global/control recursions]
\label{lem:exact-recursion}
For one communication round,
\begin{align}
  \theta_{c,h+1}
  &=\theta_{c,h}-\gamma\{f'(\theta_{c,h})+\xi_c+\eps_{c,h+1}\},
  \qquad \theta_{c,0}=x,                                           \label{eq:E1}\\
  x^+&=\frac1N\sum_{c=1}^N\theta_{c,H},                            \label{eq:E2}\\
  \xi_c^+&=\xi_c+\frac1{\gamma H}(\theta_{c,H}-x^+).              \label{eq:E3}
\end{align}
Moreover,
\begin{equation}
\label{eq:control-sum-conservation}
  \sum_{c=1}^N\xi_c^+=\sum_{c=1}^N\xi_c.
\end{equation}
Hence \cref{ass:zerosum} is invariant under the dynamics.
\end{lemma}

\begin{proof}
Equations~\eqref{eq:E1}--\eqref{eq:E3} are the defining one-round updates in the present notation. Summing Equation~\eqref{eq:E3} over clients and using Equation~\eqref{eq:E2},
\[
  \sum_c\xi_c^+
  =\sum_c\xi_c+\frac1{\gamma H}\left(\sum_c\theta_{c,H}-Nx^+\right)
  =\sum_c\xi_c.
\]
\end{proof}

\begin{corollary}[Exact global update and linear control cancellation]
\label{cor:global-exact}
Under \cref{ass:zerosum},
\begin{equation}
\label{eq:E5}
  x^+-x
  =-\gamma\sum_{h=0}^{H-1}\overline{f'(\theta_h)}
   -\gamma\sum_{h=0}^{H-1}\bar\eps_{h+1},
\end{equation}
where
\[
  \overline{f'(\theta_h)}:=\frac1N\sum_{c=1}^N f'(\theta_{c,h}),
  \qquad
  \bar\eps_h:=\frac1N\sum_{c=1}^N\eps_{c,h}.
\]
The cancellation of the linear control term is pathwise and uses no independence assumption.
\end{corollary}

\begin{proof}
Unrolling Equation~\eqref{eq:E1} gives
\[
  \theta_{c,H}-x
  =-\gamma\sum_{h=0}^{H-1}f'(\theta_{c,h})
   -\gamma H\xi_c
   -\gamma\sum_{h=0}^{H-1}\eps_{c,h+1}.
\]
Average over $c$ and use $N^{-1}\sum_c\xi_c=0$.
\end{proof}

\begin{corollary}[Exact stationary mean balance]
\label{cor:stationary-mean-balance}
At stationarity,
\begin{equation}
\label{eq:stationary-mean-balance}
  0=\sum_{h=0}^{H-1}\frac1N\sum_{c=1}^N\E[f'(\theta_{c,h})].
\end{equation}
\end{corollary}

\begin{proof}
Take the expectation in Equation~\eqref{eq:E5}. Stationarity gives $\E[x^+]=\E[x]$, while fresh mean-zero noise gives $\E[\bar\eps_h]=0$.
\end{proof}

\begin{lemma}[Exact new-control representation]
\label{lem:control-representation}
The new control admits the exact representation
\begin{equation}
\label{eq:C6-1}
\begin{aligned}
  \xi_c^+
  ={}&-\frac1H\sum_{h=0}^{H-1}
  \left(f'(\theta_{c,h})-\overline{f'(\theta_h)}\right)\\
  &-\frac1H\sum_{h=0}^{H-1}
  \left(\eps_{c,h+1}-\bar\eps_{h+1}\right).
\end{aligned}
\end{equation}
In particular, the old control $\xi_c$ has no direct term on the right-hand side.
\end{lemma}

\begin{proof}
Subtract Equation~\eqref{eq:E5} from the unrolled local endpoint, substitute the result into Equation~\eqref{eq:E3}, and cancel the term $\xi_c-\xi_c$.
\end{proof}

\begin{remark}[Mechanism already visible at the exact level]
The controls do not affect the global iterate through a surviving linear average. Any control contribution to the stationary mean must first modify the local trajectories $\theta_{c,h}$ and then return to the global update through the nonlinear map $f'(\theta_{c,h})$. This structural fact is exact and independent of the higher-order coefficient.
\end{remark}

\section{Control moments and coarse global second moment}
\label{sec:coarse}

\subsection{Uniform control moments}

\begin{lemma}[Uniform sixth control moment]
\label{lem:control-sixth}
There exist $\gamma_{0,H}^{(C6)}>0$ and $C_{\xi,6,H}<\infty$, independent of $N$ and $\gamma$, such that for $0<\gamma\le\gamma_{0,H}^{(C6)}$,
\begin{equation}
\label{eq:control-sixth}
  \frac1N\sum_{c=1}^N\E|\xi_c|^6\le C_{\xi,6,H}.
\end{equation}
Consequently, the client-averaged second and fourth control moments are uniformly bounded as well.
\end{lemma}

\begin{proof}
Write Equation~\eqref{eq:C6-1} as $\xi_c^+=-A_c-B_c$, where
\[
  A_c:=\frac1H\sum_{h=0}^{H-1}
  \left(f'(\theta_{c,h})-\overline{f'(\theta_h)}\right),
  \qquad
  B_c:=\frac1H\sum_{h=0}^{H-1}
  (\eps_{c,h+1}-\bar\eps_{h+1}).
\]
By convexity of $z\mapsto |z|^6$ and $|u-v|^6\le 32(|u|^6+|v|^6)$,
\[
  \frac1N\sum_c\E|A_c|^6
  \le 64\frac1H\sum_{h=0}^{H-1}
  \frac1N\sum_c\E|f'(\theta_{c,h})|^6.
\]
Since $f'(0)=0$ and $f'$ is $L$-Lipschitz,
\[
  |f'(y)|^6\le L^6|y|^6.
\]
The coarse local sixth-moment estimate of \citet{mangold2025scaffold}, with the sixth-moment noise envelope supplied by boundedness under \cref{ass:noise}, gives
\[
  \frac1N\sum_c\E|\theta_{c,h}|^6\le C_H\gamma^3.
\]
Hence $N^{-1}\sum_c\E|A_c|^6\le C_H\gamma^3$. On the other hand,
\[
  |\eps_{c,h}-\bar\eps_h|\le 2B_\eps
\]
pathwise, so $N^{-1}\sum_c\E|B_c|^6\le (2B_\eps)^6$. Therefore
\[
  \frac1N\sum_c\E|\xi_c^+|^6\le C_H
\]
for $\gamma\le1$. Stationarity transfers this bound from $\xi_c^+$ to $\xi_c$. The second and fourth moment bounds follow from H\"older/Jensen on the joint probability--client average.
\end{proof}

\begin{remark}
The lemma proves a \emph{client-averaged} control sixth-moment bound. No individual-client supremum is used later. In the homogeneous unique stationary law one may additionally recover exchangeability, but that strengthening is unnecessary here.
\end{remark}

\subsection{Restricted uniformity of the control second-moment expansion}

\begin{lemma}[Uniform restricted control second moment]
\label{lem:QH-uniform}
In the present homogeneous additive-noise setting,
\begin{equation}
\label{eq:QH-uniform}
  \QH=\frac{\sigma^2}{H}\left(1-\frac1N\right)+R_{Q,H},
  \qquad |R_{Q,H}|\le C_{Q,H}\gamma,
\end{equation}
where $C_{Q,H}$ is independent of $N$ and $\gamma$.
\end{lemma}

\begin{proof}
This is consistent with the specialization of \citet[Lemma~5.1]{mangold2025scaffold}; we record a restricted argument to make the uniformity explicit. From Equation~\eqref{eq:C6-1}, write
\[
  \xi_c^+=-G_c-E_c,
\]
with $G_c$ the averaged gradient-disagreement term and $E_c$ the averaged fresh-noise disagreement. For each local step,
\[
  \Var(\eps_{c,h}-\bar\eps_h)=\sigma^2\left(1-\frac1N\right),
\]
and different local-step noises are independent. Thus
\begin{equation}
\label{eq:Ec-var}
  \E[E_c^2]=\frac{\sigma^2}{H}\left(1-\frac1N\right).
\end{equation}
We first establish the finite-step estimate needed here without appealing forward to \cref{lem:displacement2}. Let
$\delt_{c,h}:=\theta_{c,h}-x$. From Equation~\eqref{eq:E1},
\[
  \delt_{c,h+1}=\delt_{c,h}
  -\gamma\{f'(x+\delt_{c,h})+\xi_c+\eps_{c,h+1}\},
  \qquad \delt_{c,0}=0.
\]
Since $f'(0)=0$ and $f'$ is $L$-Lipschitz, a finite-step discrete Gr\"onwall bound gives, for every $h\le H$,
\[
  |\delt_{c,h}|
  \le C_H\gamma\left(|x|+|\xi_c|+\sum_{j=1}^h|\eps_{c,j}|\right).
\]
    Averaging the square over clients and expectations, and using the earlier coarse bound
$\E[x^2]=O_H(\gamma)$, the coarse client-averaged control second moment, and bounded noise, yields
\[
  \frac1N\sum_c\E|\theta_{c,h}-x|^2
  =\frac1N\sum_c\E|\delt_{c,h}|^2
  \le C_H\gamma^2.
\]
Thus, the estimate used in this lemma has already been proved at this point; \cref{lem:displacement2} records it separately for later reuse.
Using smoothness and Jensen gives $N^{-1}\sum_c\E[G_c^2]\le C_H\gamma^2$. Hence
\[
  \left|\frac1N\sum_c\E[G_cE_c]\right|
  \le C_H\gamma,
\]
while $N^{-1}\sum_c\E[G_c^2]=O_H(\gamma^2)$. Expanding $\E[(\xi_c^+)^2]$, averaging clients, and using stationarity yields Equation~\eqref{eq:QH-uniform}. All constants are uniform in $N$ because $0\le1-1/N\le1$ and the preceding average-moment estimates are uniform.
\end{proof}

\subsection{Coarse global second moment}

\begin{lemma}[Finite-step local displacement]
\label{lem:displacement2}
For every $h\in\{0,\ldots,H\}$,
\begin{equation}
\label{eq:disp2}
  \frac1N\sum_{c=1}^N\E[(\theta_{c,h}-x)^2]\le C_H\gamma^2.
\end{equation}
\end{lemma}

\begin{proof}
Let $\delt_{c,h}:=\theta_{c,h}-x$. From Equation~\eqref{eq:E1},
\[
  \delt_{c,h+1}=\delt_{c,h}-\gamma\{f'(x+\delt_{c,h})+\xi_c+\eps_{c,h+1}\},
  \qquad \delt_{c,0}=0.
\]
Since $|f'(y)|\le L|y|$, finite-step discrete Gr\"onwall gives, for $h\le H$,
\[
  |\delt_{c,h}|\le C_H\gamma\left(|x|+|\xi_c|+\sum_{j=1}^h|\eps_{c,j}|\right).
\]
Use the earlier coarse bound $\E[x^2]=O_H(\gamma)$, the coarse control second moment, and bounded noise. No sharp coefficient from \cref{lem:QH-uniform} is needed.
\end{proof}

\begin{lemma}[Coarse global second moment]
\label{lem:G5a}
There are constants independent of $N,\gamma$ such that
\begin{equation}
\label{eq:G5a}
  M_2:=\E[x^2]\le C_H\smallscale
  =C_H\left(\frac{\gamma}{N}+\gamma^2\right).
\end{equation}
\end{lemma}

\begin{proof}
From \cref{cor:global-exact}, define the exact decomposition
\begin{equation}
\label{eq:G5a-decomp}
  x^+=\Phi_\gamma(x)+\zeta+\mathcal{T},
\end{equation}
where
\[
  \Phi_\gamma(x):=x-\gamma H f'(x),\qquad
  \zeta:=-\gamma\sum_{h=0}^{H-1}\bar\eps_{h+1},
\]
and
\[
  \mathcal{T}:=-\gamma\sum_{h=0}^{H-1}
  \left(\overline{f'(\theta_h)}-f'(x)\right).
\]
We call $\mathcal{T}$ the \emph{local-trajectory deviation term}; it is not purely a nonlinear remainder and need not vanish for a quadratic objective.

Client and local-step independence give
\begin{equation}
\label{eq:zeta-var}
  \E[\zeta^2]=\frac{\gamma^2H\sigma^2}{N},
  \qquad
  \E[\Phi_\gamma(x)\zeta]=0.
\end{equation}
By smoothness, Jensen, and \cref{lem:displacement2},
\begin{equation}
\label{eq:T2}
  \E[\mathcal{T}^2]\le C_H\gamma^4.
\end{equation}
Strong convexity and $\gamma HL\le1$ imply
\[
  \Phi_\gamma(x)^2\le (1-\mu H\gamma)x^2.
\]
Stationarity in Equation~\eqref{eq:G5a-decomp}, Young's inequality for the $\Phi_\gamma\mathcal{T}$ term, and Cauchy--Schwarz for $\zeta\mathcal{T}$ give
\[
  \frac{\mu H\gamma}{2}M_2
  \le C_H\left(\frac{\gamma^2}{N}+\gamma^3+\gamma^4\right).
\]
Here
\[
  \frac{\gamma^3}{\sqrt N}
  \le \frac12\left(\frac{\gamma^2}{N}+\gamma^4\right)
\]
controls the mixed noise/trajectory term without any relative-rate condition. Divide by $\gamma$ and take $\gamma\le1$.
\end{proof}

\begin{remark}[Coarse versus sharp noise interpretation]
The exact variance in Equation~\eqref{eq:zeta-var} is useful for the order bound, but $\zeta$ alone is not the sharp effective round noise for the stationary covariance: local noises also feed back through $\mathcal{T}$. The sharp covariance calculation therefore requires a separate linear/nonlinear decomposition in \cref{sec:sharpM2}.
\end{remark}

\section{Local second-moment expansion}
\label{sec:G3}

This section refines an earlier coarse local second-moment bound to the coefficient level. The direct fresh-noise term and the persistent control second moment enter through an exact square.

\begin{lemma}[Local second-moment expansion]
\label{lem:G3}
For $h\in\{0,\ldots,H\}$ define
\[
  \Ubar_h:=\frac1N\sum_{c=1}^N\left(\E[\theta_{c,h}^2]-M_2\right).
\]
Then
\begin{equation}
\label{eq:G3-QH}
  \Ubar_h
  =\gamma^2\left(h^2\QH+h\sigma^2\right)+O_H(\remscale),
\end{equation}
and, using \cref{lem:QH-uniform},
\begin{equation}
\label{eq:G3-final}
  \Ubar_h
  =\gamma^2\sigma^2\left(h+\frac{h^2}{H}\right)
   +O_H\left(\frac{\gamma^2}{N}+\gamma^3\right).
\end{equation}
The remainder is uniform in $N$ and $\gamma$ for fixed $H$.
\end{lemma}

\begin{proof}
Unroll Equation~\eqref{eq:E1} for $h$ steps and write the exact decomposition
\begin{equation}
\label{eq:G3-decomp}
  \theta_{c,h}=x+D_{c,h}+R_{c,h},
\end{equation}
where
\[
  D_{c,h}:=-\gamma\left(h\xi_c+\sum_{j=1}^h\eps_{c,j}\right),
  \qquad
  R_{c,h}:=-\gamma\sum_{\ell=0}^{h-1}f'(\theta_{c,\ell}).
\]
Expanding the square gives
\begin{equation}
\label{eq:G3-square}
\begin{aligned}
  \Ubar_h={}&
  2\frac1N\sum_c\E[xD_{c,h}]
  +\frac1N\sum_c\E[D_{c,h}^2]\\
  &+2\frac1N\sum_c\E[xR_{c,h}]
  +2\frac1N\sum_c\E[D_{c,h}R_{c,h}]
  +\frac1N\sum_c\E[R_{c,h}^2].
\end{aligned}
\end{equation}

The first term vanishes exactly. The control part vanishes after the client average because $\sum_c\xi_c=0$ pathwise, and the fresh-noise part vanishes because the round-start $x$ is independent of the mean-zero fresh noises.

For the square of $D_{c,h}$, set $S_{c,h}:=\sum_{j=1}^h\eps_{c,j}$. Conditional on the round-start state, $\E[S_{c,h}]=0$ and $\E[S_{c,h}^2]=h\sigma^2$. Hence
\[
  \E[\xi_cS_{c,h}]=0
\]
and therefore the client average satisfies the exact identity
\begin{equation}
\label{eq:D2-exact}
  \frac1N\sum_c\E[D_{c,h}^2]
  =\gamma^2\left(h^2\QH+h\sigma^2\right).
\end{equation}

It remains to bound the terms containing $R_{c,h}$. By $|f'(y)|\le L|y|$,
\[
  R_{c,h}^2\le \gamma^2hL^2\sum_{\ell=0}^{h-1}\theta_{c,\ell}^2.
\]
Using \cref{lem:G5a,lem:displacement2},
\[
  \frac1N\sum_c\E[\theta_{c,\ell}^2]\le C_H\smallscale,
\]
so
\begin{equation}
\label{eq:R2-G3}
  \frac1N\sum_c\E[R_{c,h}^2]
  \le C_H\gamma^2\smallscale
  =C_H\left(\frac{\gamma^3}{N}+\gamma^4\right).
\end{equation}
Cauchy--Schwarz then gives
\[
  \left|\frac1N\sum_c\E[xR_{c,h}]\right|
  \le C_H\gamma\smallscale=C_H\remscale.
\]
Likewise, Equation~\eqref{eq:D2-exact} and the coarse bound $\QH=O_H(1)$ give
\[
  \left|\frac1N\sum_c\E[D_{c,h}R_{c,h}]\right|
  \le C_H\gamma^2\sqrt{\smallscale}.
\]
Since
\[
  \gamma^2\sqrt{\smallscale}
  \le \frac{\gamma^{5/2}}{\sqrt N}+\gamma^3
  \le C\left(\frac{\gamma^2}{N}+\gamma^3\right),
\]
all remaining terms in Equation~\eqref{eq:G3-square} are $O_H(\remscale)$. This proves Equation~\eqref{eq:G3-QH}.

Finally, substitute Equation~\eqref{eq:QH-uniform}:
\[
  \gamma^2h^2\QH
  =\gamma^2\frac{h^2\sigma^2}{H}
   +O_H\left(\frac{\gamma^2}{N}+\gamma^3\right),
\]
which gives Equation~\eqref{eq:G3-final}.
\end{proof}

\begin{remark}[Canonical source split at the local-second-moment level]
The two leading terms in Equation~\eqref{eq:D2-exact} have distinct origins:
\[
  \underbrace{\gamma^2h\sigma^2}_{\text{direct fresh local noise}},
  \qquad
  \underbrace{\gamma^2h^2\QH}_{\text{round-start control fluctuation}}.
\]
The second term persists at order $N^0\gamma^2$ because \cref{eq:QH-uniform} has the $N^0$ limit $\sigma^2/H$. This is a source-level decomposition only; the bias coefficient is assembled later after all competing same-order channels have been excluded.
\end{remark}

\begin{remark}[Relation to the earlier coarse bound]
Lemma~E.3 of \citet{mangold2025scaffold} controls a local second-moment correction by a coarse $O(\gamma^2H)$ bound. Equation~\eqref{eq:G3-final} is the restricted coefficient-level refinement used in the present theorem; it does not attribute this coefficient to the earlier result.
\end{remark}

\section{Mean bootstrap and global higher moments}
\label{sec:mean-global}

\subsection{Coarse mean bootstrap without signed-moment circularity}

Define
\[
  q(y):=f'(y)-ay,
  \qquad a=f''(0),
\]
and
\[
  m_h:=\frac1N\sum_{c=1}^N\E[\theta_{c,h}],
  \qquad
  \nu_h:=\frac1N\sum_{c=1}^N\E[q(\theta_{c,h})],
  \qquad
  \bstat:=\E[x].
\]

\begin{lemma}[Mean bootstrap]
\label{lem:G8}
There exists $C_H<\infty$ such that
\begin{align}
  |\nu_h|&\le C_H\smallscale,\qquad h=0,\ldots,H,             \label{eq:nu-coarse}\\
  |\bstat|&\le C_H\smallscale,                                \label{eq:b-coarse}\\
  \sum_{h=0}^{H-1}m_h&=H\bstat+O_H(\remscale).                \label{eq:G8-sum}
\end{align}
This lemma uses only the coarse second-moment result \cref{lem:G5a} and the finite-step displacement estimate; it does not use \cref{lem:G3} or any signed third moment.
\end{lemma}

\begin{proof}
Since $q(0)=q'(0)=0$ and $q''=f'''$, the integral Taylor formula and $|f'''|\le K_3$ give the global bound
\begin{equation}
\label{eq:q-quadratic}
  |q(y)|\le \frac{K_3}{2}|y|^2.
\end{equation}
By \cref{lem:G5a,lem:displacement2},
\[
  \frac1N\sum_c\E[\theta_{c,h}^2]
  \le 2M_2+2\frac1N\sum_c\E[(\theta_{c,h}-x)^2]
  \le C_H\smallscale.
\]
Thus Equation~\eqref{eq:nu-coarse} follows directly from Equation~\eqref{eq:q-quadratic}.

Average Equation~\eqref{eq:E1} over clients and expectations. The zero-sum control condition and mean-zero fresh noise yield the exact recursion
\begin{equation}
\label{eq:mean-recursion}
  m_{h+1}=\lambda_\gamma m_h-\gamma\nu_h,
  \qquad
  \lambda_\gamma:=1-a\gamma,
  \qquad m_0=\bstat.
\end{equation}
Unrolling,
\begin{equation}
\label{eq:mean-unroll}
  m_h=\lambda_\gamma^h\bstat
  -\gamma\sum_{j=0}^{h-1}\lambda_\gamma^{h-1-j}\nu_j.
\end{equation}
At stationarity, $m_H=\E[x^+]=\E[x]=\bstat$. Hence
\begin{equation}
\label{eq:b-exact-weighted}
  \bstat=-\frac{1}{aA_H(\lambda_\gamma)}
  \sum_{j=0}^{H-1}\lambda_\gamma^{H-1-j}\nu_j,
  \qquad
  A_H(\lambda):=\sum_{k=0}^{H-1}\lambda^k.
\end{equation}
Choose the common step-size threshold so that $0\le\lambda_\gamma\le1$. Then the weights in Equation~\eqref{eq:b-exact-weighted} are nonnegative and sum to $A_H(\lambda_\gamma)$, giving
\[
  |\bstat|\le a^{-1}\max_j|\nu_j|\le C_H\smallscale.
\]
This proves Equation~\eqref{eq:b-coarse} without any signed-third-moment input.

Finally, subtract $\bstat$ from Equation~\eqref{eq:mean-unroll}. For fixed $h\le H$,
\[
  |\lambda_\gamma^h-1|\le ah\gamma,
\]
so Equations~\eqref{eq:nu-coarse} and~\eqref{eq:b-coarse} give
\[
  |m_h-\bstat|\le C_H\gamma\smallscale=C_H\remscale.
\]
Sum over fixed $H$ to obtain Equation~\eqref{eq:G8-sum}.
\end{proof}

\subsection{Exact linear/nonlinear global split}

For the remaining global-moment arguments, define
\[
  \rho_\gamma:=\lambda_\gamma^H,
\]
and the exact decomposition obtained by unrolling the local recursion written as
$f'(y)=ay+q(y)$:
\begin{equation}
\label{eq:global-linear-nonlinear}
  x^+=\rho_\gamma x+\eta+\Delta,
\end{equation}
where
\begin{align}
  \eta&:=-\gamma\sum_{j=1}^H\lambda_\gamma^{H-j}\bar\eps_j,             \label{eq:eta-def}\\
  \Delta&:=-\gamma\sum_{h=0}^{H-1}\lambda_\gamma^{H-1-h}\bar q_h,
  \qquad
  \bar q_h:=\frac1N\sum_{c=1}^N q(\theta_{c,h}).                         \label{eq:Delta-def}
\end{align}
The linear control term cancels exactly after the client average.

The fresh-noise term satisfies
\begin{align}
  \E[\eta]&=0,                                                           \label{eq:eta-mean}\\
  \E[\eta^2]&=\frac{\gamma^2\sigma^2}{N}
  \sum_{k=0}^{H-1}\lambda_\gamma^{2k},                                  \label{eq:eta2}\\
  |\E[\eta^3]|&\le C_H\frac{\gamma^3}{N^2},                            \label{eq:eta3}\\
  \E[\eta^4]&\le C_H\frac{\gamma^4}{N^2}.                             \label{eq:eta4}
\end{align}
The third-moment estimate holds without noise symmetry.

\begin{lemma}[Nonlinear-increment moments]
\label{lem:Delta-moments}
For the $\Delta$ in Equation~\eqref{eq:Delta-def},
\begin{align}
  \E[\Delta^2]&\le C_H\gamma^2\smallscale,                              \label{eq:Delta2}\\
  \E[\Delta^4]&\le C_H\gamma^7,                                         \label{eq:Delta4-sharp}\\
  \E|\Delta|^6&\le C_H\gamma^9.                                         \label{eq:Delta6}
\end{align}
A second, weaker but sometimes useful estimate is
\begin{equation}
\label{eq:Delta4-M4}
  \E[\Delta^4]\le C_H\gamma^4(M_4+\gamma^4),
  \qquad M_4:=\E[x^4].
\end{equation}
\end{lemma}

\begin{proof}
Besides the quadratic bound in Equation~\eqref{eq:q-quadratic}, smoothness gives the global linear bound
\[
  |q(y)|\le (L+a)|y|.
\]
The earlier coarse local fourth- and sixth-moment bounds give, uniformly for $h\le H$,
\[
  \frac1N\sum_c\E|\theta_{c,h}|^4\le C_H\gamma^2,
  \qquad
  \frac1N\sum_c\E|\theta_{c,h}|^6\le C_H\gamma^3.
\]
Using the quadratic bound, Jensen, and fixed $H$ yields $\E[\Delta^2]\le C_H\gamma^4$, which is stronger than Equation~\eqref{eq:Delta2} because $\smallscale\ge\gamma^2$.

For the fourth moment, combine the linear and quadratic bounds to obtain
\[
  |q(y)|^4\le C|y|^6.
\]
Jensen then yields
\[
  \E[\Delta^4]\le C_H\gamma^4
  \sum_{h=0}^{H-1}\frac1N\sum_c\E|\theta_{c,h}|^6
  \le C_H\gamma^7,
\]
which is the sharpened estimate needed for the fourth-moment absorption argument.

For completeness, we now prove the weaker estimate in Equation~\eqref{eq:Delta4-M4} without appealing to a later lemma. The same finite-step argument used in \cref{lem:displacement2} gives pathwise, for $\delt_{c,h}:=\theta_{c,h}-x$ and fixed $h\le H$,
\[
  |\delt_{c,h}|
  \le C_H\gamma\left(|x|+|\xi_c|+\sum_{j=1}^h|\eps_{c,j}|\right).
\]
Raise this inequality to the fourth power, average over clients and expectations, and use the earlier coarse global fourth-moment bound, the client-averaged fourth control moment from \cref{lem:control-sixth}, and bounded noise. This yields
\[
  \frac1N\sum_c\E|\delt_{c,h}|^4\le C_H\gamma^4.
\]
Consequently,
\[
  \frac1N\sum_c\E|\theta_{c,h}|^4
  \le 8M_4+8\frac1N\sum_c\E|\delt_{c,h}|^4
  \le C_H(M_4+\gamma^4).
\]
Using the global linear bound $|q(y)|\le (L+a)|y|$, Jensen over clients and the fixed sum over local steps now gives
\[
  \E[\Delta^4]
  \le C_H\gamma^4\sum_{h=0}^{H-1}\frac1N\sum_c\E|\theta_{c,h}|^4
  \le C_H\gamma^4(M_4+\gamma^4),
\]
which is exactly Equation~\eqref{eq:Delta4-M4}. Thus no forward reference to the later displacement-fourth-moment lemma is needed.

Finally, $|q(y)|^6\le C|y|^6$ by the linear bound, so the same argument gives Equation~\eqref{eq:Delta6}.
\end{proof}

\subsection{Global fourth moment}

\begin{lemma}[Global fourth moment]
\label{lem:M4}
There exists $C_H<\infty$ such that
\begin{equation}
\label{eq:M4}
  M_4:=\E[x^4]\le C_H\fourthscale
  =C_H\left(\frac{\gamma^2}{N^2}+\gamma^3\right).
\end{equation}
\end{lemma}

\begin{proof}
Let $Y:=\rho_\gamma x+\eta$. Stationarity and Equation~\eqref{eq:global-linear-nonlinear} give
\begin{equation}
\label{eq:M4-stationary}
\begin{aligned}
  M_4={}&\E[Y^4]
  +4\E[Y^3\Delta]
  +6\E[Y^2\Delta^2]
  +4\E[Y\Delta^3]
  +\E[\Delta^4].
\end{aligned}
\end{equation}
Because $\eta$ is independent of the round-start $x$ and has zero mean,
\begin{equation}
\label{eq:Y4-expand}
  \E[Y^4]
  =\rho_\gamma^4M_4
   +6\rho_\gamma^2M_2\E[\eta^2]
   +4\rho_\gamma\bstat\E[\eta^3]
   +\E[\eta^4].
\end{equation}
By \cref{lem:G5a}, Equations~\eqref{eq:eta2}--\eqref{eq:eta4}, and $|\bstat|\le\sqrt{M_2}$,
\[
  M_2\E[\eta^2]+|\bstat\E[\eta^3]|+\E[\eta^4]
  \le C_H\gamma\fourthscale.
\]

For the nonlinear terms, fix $\epsilon>0$. Young's inequality and the sharpened Equation~\eqref{eq:Delta4-sharp} give
\begin{align*}
  |\E[Y^3\Delta]|
  &\le \epsilon\gamma\E[Y^4]+C_{H,\epsilon}\gamma^{-3}\E[\Delta^4]
   \le C_H\epsilon\gamma M_4+C_{H,\epsilon}\gamma\fourthscale,\\
  |\E[Y^2\Delta^2]|
  &\le \epsilon\gamma\E[Y^4]+C_{H,\epsilon}\gamma^{-1}\E[\Delta^4]
   \le C_H\epsilon\gamma M_4+C_{H,\epsilon}\gamma\fourthscale,\\
  |\E[Y\Delta^3]|
  &\le \epsilon\gamma\E[Y^4]+C_{H,\epsilon}\gamma^{-1/3}\E[\Delta^4]
   \le C_H\epsilon\gamma M_4+C_{H,\epsilon}\gamma\fourthscale,
\end{align*}
and $\E[\Delta^4]\le C_H\gamma\fourthscale$ for $\gamma\le1$.

Finally,
\[
  1-\rho_\gamma^4
  =a\gamma\sum_{k=0}^{4H-1}(1-a\gamma)^k
  \ge a\gamma
\]
when $0\le1-a\gamma\le1$. Move the $\rho_\gamma^4M_4$ term to the left of Equation~\eqref{eq:M4-stationary}, choose a fixed $\epsilon$ small enough that the $\epsilon\gamma M_4$ terms are absorbed, and divide by $\gamma$ to obtain Equation~\eqref{eq:M4}.
\end{proof}

\begin{proofnote}[Why Equation~\eqref{eq:Delta4-sharp} is needed]
The weaker estimate in Equation~\eqref{eq:Delta4-M4} is true but does not by itself justify the $Y^3\Delta$ absorption: after the factor $\gamma^{-3}$ from Young's inequality, it produces a coefficient of order $C_\epsilon\gamma M_4$, whose constant worsens as $\epsilon\downarrow0$. The independent bound $\E[\Delta^4]\le C_H\gamma^7$ removes this defect. The sharper bound is needed to complete the absorption argument.
\end{proofnote}

\subsection{Global signed third moment}

\begin{lemma}[Global signed third moment]
\label{lem:M3}
There exists $C_H<\infty$ such that
\begin{equation}
\label{eq:M3-target}
  |M_3|:=|\E[x^3]|\le C_H\remscale.
\end{equation}
In fact, the proof below yields the stronger internal estimate
\begin{equation}
\label{eq:M3-strong}
  |M_3|\le C_H\fourthscale.
\end{equation}
Only Equation~\eqref{eq:M3-target} is needed for the main theorem.
\end{lemma}

\begin{proof}
Again set $Y:=\rho_\gamma x+\eta$. Stationarity gives the exact identity
\begin{equation}
\label{eq:M3-stationary}
\begin{aligned}
  (1-\rho_\gamma^3)M_3={}&
  3\rho_\gamma\bstat\E[\eta^2]+\E[\eta^3]\\
  &+3\E[Y^2\Delta]
  +3\E[Y\Delta^2]
  +\E[\Delta^3].
\end{aligned}
\end{equation}
The first two terms are $O_H(\gamma\fourthscale)$ by \cref{lem:G8} and Equations~\eqref{eq:eta2} and~\eqref{eq:eta3}.

The delicate term is $\E[Y^2\Delta]$. A generic Cauchy--Schwarz bound is too coarse, so use the structure in Equations~\eqref{eq:Delta-def} and~\eqref{eq:q-quadratic}:
\[
  |\E[Y^2\Delta]|
  \le C_H\gamma\sum_{h=0}^{H-1}
  \frac1N\sum_c\E[Y^2\theta_{c,h}^2].
\]
Let $\delt_{c,h}:=\theta_{c,h}-x$. A finite-step fourth-moment argument using \cref{lem:control-sixth,lem:M4} gives
\begin{equation}
\label{eq:disp4}
  \frac1N\sum_c\E[\delt_{c,h}^4]\le C_H\gamma^4.
\end{equation}
Since $Y^2\le C_H(x^2+\eta^2)$ and $\theta_{c,h}^2\le2x^2+2\delt_{c,h}^2$, the joint probability--client average is bounded by a constant times
\[
  M_4
  +\frac1N\sum_c\E[x^2\delt_{c,h}^2]
  +\E[\eta^2x^2]
  +\frac1N\sum_c\E[\eta^2\delt_{c,h}^2].
\]
The first three terms are
$O_H(\fourthscale)$,
$O_H(\sqrt{\fourthscale\,\gamma^4})=O_H(\fourthscale)$, and
$O_H(\E[\eta^2]M_2)=O_H(\fourthscale)$.
The fourth is
\[
  O_H\!\left(\sqrt{\E[\eta^4]\,\frac1N\sum_c\E[\delt_{c,h}^4]}\right)
  =O_H\left(\frac{\gamma^4}{N}\right)
  =O_H(\fourthscale).
\]
Thus
\begin{equation}
\label{eq:Y2Delta}
  |\E[Y^2\Delta]|\le C_H\gamma\fourthscale.
\end{equation}

Next,
\[
  |\E[Y\Delta^2]|
  \le \sqrt{\E[Y^2]\E[\Delta^4]}
  \le C_H\gamma^4
  \le C_H\gamma\fourthscale,
\]
where \cref{lem:G5a,lem:Delta-moments} were used. Finally,
\[
  |\E[\Delta^3]|\le \sqrt{\E|\Delta|^6}\le C_H\gamma^{9/2}
  \le C_H\gamma\fourthscale.
\]
The restoring factor satisfies
\[
  1-\rho_\gamma^3
  =a\gamma\sum_{k=0}^{3H-1}(1-a\gamma)^k
  \ge a\gamma.
\]
Divide Equation~\eqref{eq:M3-stationary} by this factor to obtain Equation~\eqref{eq:M3-strong}, and hence Equation~\eqref{eq:M3-target}. Notice that the skewness contribution $\E[\eta^3]$ is retained explicitly; symmetry is unnecessary.
\end{proof}

\section{Sharp global second moment}
\label{sec:sharpM2}

This is the key identifiability step. It rules out an unresolved client-number-independent $N^0\gamma^2$ contribution from the global raw second moment itself.

\begin{lemma}[Averaged local displacement]
\label{lem:avg-delta}
Let
\[
  \bar\delt_h:=\frac1N\sum_{c=1}^N(\theta_{c,h}-x).
\]
Then, for fixed $h\le H$,
\begin{equation}
\label{eq:avg-delta2}
  \E[\bar\delt_h^2]
  \le C_H\gamma^2\left(\smallscale+\frac1N\right).
\end{equation}
\end{lemma}

\begin{proof}
Average the local displacement recursion. The control term cancels pathwise, so
\[
  \bar\delt_{h+1}
  =\bar\delt_h-\gamma\overline{f'(\theta_h)}-\gamma\bar\eps_{h+1},
  \qquad \bar\delt_0=0.
\]
Unrolling,
\[
  \bar\delt_h=-\gamma\sum_{\ell=0}^{h-1}\overline{f'(\theta_\ell)}
  -\gamma\sum_{j=1}^h\bar\eps_j.
\]
For fixed $H$, Jensen's inequality, smoothness, and the coarse local second-moment estimate give
\[
  \E|\overline{f'(\theta_\ell)}|^2
  \le L^2\frac1N\sum_c\E[\theta_{c,\ell}^2]
  \le C_H\smallscale,
\]
while $\E[\bar\eps_j^2]=\sigma^2/N$. This gives Equation~\eqref{eq:avg-delta2}.
\end{proof}

\begin{lemma}[Sharp global second moment]
\label{lem:G5b}
There exists $C_H<\infty$ such that
\begin{equation}
\label{eq:G5b}
  M_2
  =\frac{\sigma^2}{2a}\frac{\gamma}{N}
  +O_H\left(\frac{\gamma^2}{N}+\gamma^3\right).
\end{equation}
Equivalently, the nonlinear correction to the stationary second moment has no unresolved $N^0\gamma^2$ term.
\end{lemma}

\begin{proof}
We use the exact decomposition in Equation~\eqref{eq:global-linear-nonlinear}.

\paragraph{Step 1: exact linearized stationary second moment.}
For the linearized chain $x_{\mathrm{lin}}^+=\rho_\gamma x_{\mathrm{lin}}+\eta$, stationarity gives
\[
  M_{2,\mathrm{lin}}=\frac{\E[\eta^2]}{1-\rho_\gamma^2}.
\]
By Equation~\eqref{eq:eta2},
\[
  \E[\eta^2]
  =\frac{\gamma^2\sigma^2}{N}\sum_{k=0}^{H-1}\lambda_\gamma^{2k}.
\]
Since
\[
  1-\rho_\gamma^2
  =1-\lambda_\gamma^{2H}
  =(1-\lambda_\gamma^2)\sum_{k=0}^{H-1}\lambda_\gamma^{2k},
\]
the $H$-dependent geometric sum cancels exactly:
\begin{equation}
\label{eq:M2lin-exact}
  M_{2,\mathrm{lin}}
  =\frac{\gamma\sigma^2}{aN(2-a\gamma)}.
\end{equation}
Therefore
\begin{equation}
\label{eq:M2lin-first}
  M_{2,\mathrm{lin}}
  =\frac{\sigma^2}{2a}\frac{\gamma}{N}
  +O\left(\frac{\gamma^2}{N}\right).
\end{equation}
The explicit $+\sigma^2\gamma^2/(4N)$ term in the Taylor expansion of Equation~\eqref{eq:M2lin-exact} is \emph{not} claimed to be the complete nonlinear finite-$N$ coefficient.

\paragraph{Step 2: exact nonlinear stationary equation.}
Squaring Equation~\eqref{eq:global-linear-nonlinear} and using $\E[x\eta]=0$ gives
\begin{equation}
\label{eq:M2-nonlinear-eq}
  (1-\rho_\gamma^2)M_2
  =\E[\eta^2]
  +2\rho_\gamma\E[x\Delta]
  +2\E[\eta\Delta]
  +\E[\Delta^2].
\end{equation}
By \cref{lem:Delta-moments},
\begin{equation}
\label{eq:Delta2-rscale}
  \E[\Delta^2]\le C_H\gamma\remscale.
\end{equation}
It remains to obtain equally sharp bounds for the two covariance terms.

\paragraph{Step 3: the signed covariance $\E[x\Delta]$.}
For $\delt_{c,h}:=\theta_{c,h}-x$, Taylor-expand $q$ around the round-start $x$:
\begin{equation}
\label{eq:q-around-x}
  q(x+\delt_{c,h})
  =q(x)+q'(x)\delt_{c,h}+r_{c,h},
  \qquad
  |r_{c,h}|\le \frac{K_3}{2}\delt_{c,h}^2.
\end{equation}
After the client average, we write
\[
  \bar q_h=q(x)+q'(x)\bar\delt_h+\bar r_h
\]
and hence
\begin{equation}
\label{eq:Delta-split}
  \Delta=\Delta_0+R_1+R_2,
\end{equation}
with
\begin{align*}
  \Delta_0&:=-\gamma A_H(\lambda_\gamma)q(x),\\
  R_1&:=-\gamma\sum_{h=0}^{H-1}\lambda_\gamma^{H-1-h}q'(x)\bar\delt_h,\\
  R_2&:=-\gamma\sum_{h=0}^{H-1}\lambda_\gamma^{H-1-h}\bar r_h.
\end{align*}
Since $q(y)=\frac\tau2y^2+O(|y|^3)$, we have
\[
  |\E[xq(x)]|\le C(|M_3|+M_4)
  \le C_H\remscale
\]
by \cref{lem:M3,lem:M4}. Thus
\begin{equation}
\label{eq:xDelta0}
  |\E[x\Delta_0]|\le C_H\gamma\remscale.
\end{equation}
Moreover $|q'(x)|\le K_3|x|$, so \cref{lem:M4,lem:avg-delta} give
\[
  |\E[xR_1]|
  \le C_H\gamma\sqrt{M_4\,\E[\bar\delt_h^2]}
  \le C_H\gamma^2\sqrt{\fourthscale\left(\smallscale+\frac1N\right)}.
\]
Because $\fourthscale\le\gamma\smallscale$ and
$\gamma(\smallscale+1/N)\le 2\smallscale$ for $\gamma\le1$,
\begin{equation}
\label{eq:xR1}
  |\E[xR_1]|\le C_H\gamma^2\smallscale=C_H\gamma\remscale.
\end{equation}
For $R_2$, Equation~\eqref{eq:disp4} and Jensen imply $\E[\bar r_h^2]\le C_H\gamma^4$, hence
\[
  |\E[xR_2]|
  \le C_H\gamma\sqrt{M_2\gamma^4}
  =C_H\gamma^3\sqrt{\smallscale}
  \le C_H\gamma^2\smallscale.
\]
Combining the three pieces,
\begin{equation}
\label{eq:xDelta-sharp}
  |\E[x\Delta]|\le C_H\gamma\remscale.
\end{equation}

\paragraph{Step 4: the fresh-noise covariance $\E[\eta\Delta]$ by coordinate replacement.}
A naive Cauchy--Schwarz bound is too coarse. First note that $\Delta_0$ in Equation~\eqref{eq:Delta-split} depends only on the round-start $x$, so
\begin{equation}
\label{eq:etaDelta0}
  \E[\eta\Delta_0]=0.
\end{equation}
Let $R:=\Delta-\Delta_0$. Index the $NH$ fresh noise coordinates by $i=(c,j)$ and write
\[
  \eta=\sum_i\alpha_i\eps_i,
  \qquad |\alpha_i|\le \frac{\gamma}{N}.
\]
For a fixed coordinate $i$, let $Z^{(i,0)}$ denote the noise array obtained by replacing $\eps_i$ by $0$ while keeping all other coordinates unchanged. Since $R(Z^{(i,0)})$ is independent of $\eps_i$ and $\E[\eps_i]=0$,
\begin{equation}
\label{eq:coord-replace}
  \E[\eps_iR(Z)]
  =\E\left[\eps_i\{R(Z)-R(Z^{(i,0)})\}\right].
\end{equation}
We now justify the required path sensitivity deterministically. Write $i=(c_0,j_0)$ and couple the two within-round trajectories by using the same round-start state and the same fresh noises except that the coordinate $\eps_{c_0,j_0}$ is replaced by $0$ in $Z^{(i,0)}$. All clients $c\neq c_0$ then have identical trajectories. For the affected client, we define
\[
  d_h:=\theta_{c_0,h}(Z)-\theta_{c_0,h}(Z^{(i,0)}).
\]
Before the perturbed noise is used, $d_h=0$ for $h<j_0$, while the update containing that coordinate gives
\[
  d_{j_0}=-\gamma\eps_i.
\]
For every subsequent local step $h\ge j_0$, the two recursions use the same control and the same remaining fresh noises, hence
\[
  d_{h+1}
  =d_h-\gamma\{f'(\theta_{c_0,h}(Z))-f'(\theta_{c_0,h}(Z^{(i,0)}))\}.
\]
By $L$-smoothness,
\[
  |d_{h+1}|\le (1+\gamma L)|d_h|,
\]
and therefore, for $j_0\le h\le H$,
\[
  |d_h|
  \le \gamma(1+\gamma L)^{h-j_0}|\eps_i|
  \le C_H\gamma|\eps_i|,
\]
where the last constant is uniform in $N$ and $\gamma$ under the common small-step restriction.

Because $\Delta_0$ depends only on the round-start $x$, the difference $R=\Delta-\Delta_0$ is the same as the difference of $\Delta$. Only client $c_0$ contributes to the change in each client average $\bar q_h$. Since
\[
  q'(y)=f''(y)-a,
  \qquad |q'(y)|\le L+a,
\]
we obtain
\[
  |\bar q_h(Z)-\bar q_h(Z^{(i,0)})|
  \le \frac{L+a}{N}|d_h|
  \le C_H\frac{\gamma}{N}|\eps_i|.
\]
Moreover, under the same threshold $0\le\lambda_\gamma\le1$, the weights in Equation~\eqref{eq:Delta-def} have absolute value at most one. Summing at most $H$ affected local steps and using the outer factor $\gamma$ in Equation~\eqref{eq:Delta-def} yields the deterministic sensitivity bound
\begin{equation}
\label{eq:R-coordinate-Lipschitz}
  |R(Z)-R(Z^{(i,0)})|
  \le C_H\frac{\gamma^2}{N}|\eps_i|.
\end{equation}
Hence
\[
  |\E[\eps_iR]|\le C_H\frac{\gamma^2}{N}\sigma^2.
\]
Summing the $NH$ coordinates in $\eta$,
\begin{equation}
\label{eq:etaDelta-sharp}
  |\E[\eta\Delta]|
  =|\E[\eta R]|
  \le NH\left(C_H\frac{\gamma}{N}\right)
          \left(C_H\frac{\gamma^2}{N}\right)
  \le C_H\frac{\gamma^3}{N}.
\end{equation}
This is the extra $1/N$ gain that the generic Cauchy--Schwarz route misses.

\paragraph{Step 5: conclude.}
Insert Equations~\eqref{eq:Delta2-rscale}, \eqref{eq:xDelta-sharp}, and~\eqref{eq:etaDelta-sharp} into Equation~\eqref{eq:M2-nonlinear-eq}:
\[
  (1-\rho_\gamma^2)M_2
  =\E[\eta^2]+O_H\left(\frac{\gamma^3}{N}+\gamma^4\right).
\]
Because
\[
  1-\rho_\gamma^2
  =a\gamma\sum_{k=0}^{2H-1}(1-a\gamma)^k
  \ge a\gamma,
\]
we obtain
\[
  M_2=M_{2,\mathrm{lin}}+O_H\left(\frac{\gamma^2}{N}+\gamma^3\right).
\]
Use Equation~\eqref{eq:M2lin-first} to conclude Equation~\eqref{eq:G5b}.
\end{proof}

\begin{remark}[What \cref{lem:G5b} does and does not identify]
The lemma excludes an unresolved client-number-independent $N^0\gamma^2$ term in $M_2$. It does \emph{not} identify the complete coefficient of the finite-$N$ monomial $\gamma^2/N$: the nonlinear correction is itself only controlled at $O_H(\gamma^2/N)$. Accordingly, the linearized expansion's coefficient $+\sigma^2/4$ at $\gamma^2/N$ must not be promoted to a theorem about the full nonlinear process.
\end{remark}

\section{Local higher moments and Taylor control}
\label{sec:local-higher}

This section places the cubic and quartic pieces of the gradient expansion inside the uniform remainder, thereby excluding an additional client-number-independent $N^0\gamma^2$ contribution from those terms.

\begin{lemma}[Local displacement fourth moment]
\label{lem:disp4}
For every fixed $h\le H$,
\begin{equation}
\label{eq:disp4-local}
  \frac1N\sum_{c=1}^N\E|\theta_{c,h}-x|^4\le C_H\gamma^4.
\end{equation}
\end{lemma}

\begin{proof}
The finite-step bound used in \cref{lem:displacement2} also yields
\[
  |\theta_{c,h}-x|
  \le C_H\gamma\left(|x|+|\xi_c|+\sum_{j=1}^h|\eps_{c,j}|\right).
\]
Raise to the fourth power, average over clients, and use \cref{lem:M4,lem:control-sixth} together with bounded noise. The resulting constant is independent of $N$ and $\gamma$ for fixed $H$.
\end{proof}

\begin{lemma}[Local signed cubic and fourth moments]
\label{lem:G6local}
For every fixed $h\le H$,
\begin{align}
  \left|\frac1N\sum_{c=1}^N\E[\theta_{c,h}^3]\right|
  &\le C_H\remscale,                                                     \label{eq:G6-cubic}\\
  \frac1N\sum_{c=1}^N\E[\theta_{c,h}^4]
  &\le C_H\fourthscale.                                                  \label{eq:G6-fourth}
\end{align}
The proof uses the coarse $M_2$ bound, the global $M_3/M_4$ bounds, and control/displacement moments; it does not require the sharp coefficient in \cref{lem:G5b}.
\end{lemma}

\begin{proof}
Write $\delt_{c,h}:=\theta_{c,h}-x$ and $\bar\delt_h=N^{-1}\sum_c\delt_{c,h}$. We already have
\[
  \frac1N\sum_c\E[\delt_{c,h}^2]\le C_H\gamma^2,
  \qquad
  \frac1N\sum_c\E[\delt_{c,h}^4]\le C_H\gamma^4,
\]
from \cref{lem:displacement2,lem:disp4}, and
\[
  \E[\bar\delt_h^2]\le C_H\gamma^2\left(\smallscale+\frac1N\right)
\]
from \cref{lem:avg-delta}.

For the signed cubic, expand exactly:
\begin{equation}
\label{eq:local-cubic-expand}
\begin{aligned}
  \frac1N\sum_c\E[\theta_{c,h}^3]
  ={}&M_3
  +3\E[x^2\bar\delt_h]
  +3\E\left[x\frac1N\sum_c\delt_{c,h}^2\right]\\
  &+\E\left[\frac1N\sum_c\delt_{c,h}^3\right].
\end{aligned}
\end{equation}
The first term is $O_H(\remscale)$ by \cref{lem:M3}. For the second term,
\[
  |\E[x^2\bar\delt_h]|
  \le \sqrt{M_4\E[\bar\delt_h^2]}
  \le C_H\remscale;
\]
indeed $\fourthscale=\gamma^2(N^{-2}+\gamma)$ and
$\smallscale+N^{-1}\le C(N^{-1}+\gamma)$ for $\gamma\le1$.
For the third term, Jensen and Cauchy--Schwarz give
\[
  \left|\E\left[x\frac1N\sum_c\delt_{c,h}^2\right]\right|
  \le \sqrt{M_2\frac1N\sum_c\E[\delt_{c,h}^4]}
  \le C_H\gamma^2\sqrt{\smallscale}
  \le C_H\remscale.
\]
Finally,
\[
  \left|\E\left[\frac1N\sum_c\delt_{c,h}^3\right]\right|
  \le \left(\frac1N\sum_c\E\delt_{c,h}^2\right)^{1/2}
       \left(\frac1N\sum_c\E\delt_{c,h}^4\right)^{1/2}
  \le C_H\gamma^3
  \le C_H\remscale.
\]
This proves Equation~\eqref{eq:G6-cubic}. Notice that only the \emph{signed client average} is controlled; no bound of the stronger form $N^{-1}\sum_c|\E[\theta_{c,h}^3]|$ is claimed.

For the fourth moment,
\[
  |x+\delt|^4\le 8(|x|^4+|\delt|^4),
\]
so
\[
  \frac1N\sum_c\E[\theta_{c,h}^4]
  \le 8M_4+8\frac1N\sum_c\E[\delt_{c,h}^4]
  \le C_H(\fourthscale+\gamma^4)
  \le C_H\fourthscale.
\]
\end{proof}

\begin{remark}[No symmetry requirement]
The proof of \cref{lem:G6local} uses mean-zero fresh noise and the global signed-third estimate from \cref{lem:M3}; it does not require $\E[\eps^3]=0$. The asymmetric-noise skew term was already retained and controlled in \cref{lem:M3}.
\end{remark}

\subsection{Gradient Taylor remainder}

\begin{lemma}[General $C^5$ Taylor control]
\label{lem:G7}
For
\[
  a=f''(0),\qquad \tau=f'''(0),\qquad \kappa=f''''(0),
\]
the gradient admits
\begin{equation}
\label{eq:grad-taylor}
  f'(y)=ay+\frac\tau2y^2+\frac\kappa6y^3+\rho_5(y),
  \qquad
  |\rho_5(y)|\le \frac{K_5}{24}|y|^4.
\end{equation}
Moreover,
\begin{equation}
\label{eq:G7}
  \left|
  \sum_{h=0}^{H-1}\frac1N\sum_{c=1}^N
  \E\left[
  f'(\theta_{c,h})-a\theta_{c,h}-\frac\tau2\theta_{c,h}^2
  \right]
  \right|
  \le C_H\remscale.
\end{equation}
\end{lemma}

\begin{proof}
Apply Taylor's theorem to $g=f'$ through cubic order. Since $g^{(4)}=f^{(5)}$,
\[
  \rho_5(y)=\frac{y^4}{3!}\int_0^1(1-s)^3 f^{(5)}(sy)\,ds,
\]
which gives the coefficient $K_5/24$ in Equation~\eqref{eq:grad-taylor}.

For each $h$, subtract the linear and quadratic pieces:
\[
  \frac1N\sum_c\E\left[
  f'(\theta_{c,h})-a\theta_{c,h}-\frac\tau2\theta_{c,h}^2
  \right]
  =\frac\kappa6\frac1N\sum_c\E[\theta_{c,h}^3]
   +\frac1N\sum_c\E[\rho_5(\theta_{c,h})].
\]
By \cref{lem:G6local}, the first term is $O_H(\remscale)$ and the second is $O_H(\fourthscale)$. Since $\fourthscale\le\remscale$ for $N\ge2$, each local-step contribution is $O_H(\remscale)$, and summation over fixed $H$ preserves this order.
\end{proof}

\begin{remark}[Precise interpretation of the higher-curvature terms]
The cubic-gradient term proportional to $f''''(0)$ and the quartic remainder of the gradient expansion, controlled by $f^{(5)}$, do not alter the client-number-independent $N^0\gamma^2$ coefficient. They may still contribute to finite-$N$ terms of order $\gamma^2/N$; no complete fixed-$N$ second-order coefficient is claimed.
\end{remark}

\begin{remark}[Why homogeneity matters for the signed cubic]
Because all clients share the same coefficient $\kappa=f''''(0)$, the Taylor contribution is proportional to the signed average $N^{-1}\sum_c\E[\theta_{c,h}^3]$. With heterogeneous clients, different fourth derivatives would weight the client cubics differently, and the present signed-average estimate would not be sufficient without additional control.
\end{remark}

\section{Assembly proof of the main expansion}
\label{app:main-proof}

This appendix proves \cref{thm:main} from the preceding lemmas.

\begin{figure}[H]
\centering
\begin{tikzpicture}[
  node distance=4mm,
  box/.style={draw,rounded corners,align=center,inner sep=4pt,text width=0.72\linewidth},
  arr/.style={-{Latex[length=2mm]},thick}
]
\node[box] (a) {Stationarity inputs + exact \mbox{\Scaf{}} recursion\\+ zero-sum control state};
\node[box,below=of a] (b) {Uniform control moments and coarse global/local second moments};
\node[box,below=of b] (c) {Coefficient-level local second-moment expansion};
\node[box,below=of c] (d) {Mean bootstrap + global fourth moment\\+ signed global third moment};
\node[box,below=of d] (e) {Sharp global second moment: no unresolved client-independent $\gamma^2$ term in $\E[x^2]$};
\node[box,below=of e] (f) {Local higher moments + $C^5$ Taylor remainder control};
\node[box,below=of f] (g) {Stationary mean assembly and uniform joint remainder};
\draw[arr] (a) -- (b);
\draw[arr] (b) -- (c);
\draw[arr] (c) -- (d);
\draw[arr] (d) -- (e);
\draw[arr] (e) -- (f);
\draw[arr] (f) -- (g);
\end{tikzpicture}
\caption{Proof dependency map. Each stage either identifies a source term or excludes a competing client-independent contribution at order $\gamma^2$. The detailed lemmas and constants are given in the surrounding appendices.}
\label{fig:proof-map}
\end{figure}

\begin{proof}[Proof of \cref{thm:main}]
The exact stationary mean balance in Equation~\eqref{eq:stationary-mean-balance} and \cref{lem:G7} give
\begin{equation}
\label{eq:assembly1}
  0
  =a\sum_{h=0}^{H-1}m_h
   +\frac\tau2\sum_{h=0}^{H-1}s_h
   +O_H(\remscale),
\end{equation}
where
\[
  s_h:=\frac1N\sum_c\E[\theta_{c,h}^2].
\]
By \cref{lem:G8},
\[
  \sum_{h=0}^{H-1}m_h=H\bstat+O_H(\remscale).
\]
By definition of \cref{lem:G3},
\[
  s_h=M_2+\Ubar_h.
\]
Substituting into Equation~\eqref{eq:assembly1},
\[
  aH\bstat+\frac\tau2HM_2
  +\frac\tau2\sum_{h=0}^{H-1}\Ubar_h
  +O_H(\remscale)=0.
\]
Hence
\begin{equation}
\label{eq:b-assembly}
  \bstat
  =-\frac\tau{2a}M_2
   -\frac\tau{2aH}\sum_{h=0}^{H-1}\Ubar_h
   +O_H(\remscale).
\end{equation}
The sharp global second moment \cref{lem:G5b} yields
\begin{equation}
\label{eq:assembly-M2}
  -\frac\tau{2a}M_2
  =-\frac{\tau\sigma^2}{4a^2}\frac{\gamma}{N}
   +O_H(\remscale).
\end{equation}
For the local correction, \cref{lem:G3} gives
\[
  \sum_{h=0}^{H-1}\Ubar_h
  =\gamma^2\sigma^2\left(
    \sum_{h=0}^{H-1}h+
    \frac1H\sum_{h=0}^{H-1}h^2
  \right)+O_H(\remscale).
\]
Use
\[
  \sum_{h=0}^{H-1}h=\frac{H(H-1)}2,
  \qquad
  \sum_{h=0}^{H-1}h^2=\frac{H(H-1)(2H-1)}6.
\]
Then
\[
  \sum_{h=0}^{H-1}h+\frac1H\sum_{h=0}^{H-1}h^2
  =\frac{(H-1)(5H-1)}6,
\]
and therefore
\begin{equation}
\label{eq:assembly-U}
  -\frac\tau{2aH}\sum_{h=0}^{H-1}\Ubar_h
  =-\frac{\tau\sigma^2}{12a}
  \frac{(H-1)(5H-1)}{H}\gamma^2
  +O_H(\remscale).
\end{equation}
Combining Equations~\eqref{eq:b-assembly}--\eqref{eq:assembly-U} gives Equation~\eqref{eq:main-expansion}. The common small-step threshold and $N,\gamma$-uniform remainder constant are justified in Appendix~\ref{sec:uniformity}.
\end{proof}

\subsection{Canonical source-wise decomposition}

\begin{corollary}[Direct local-noise and SCAFFOLD control-second-moment sources]
\label{cor:decomp}
The $N^0\gamma^2$ coefficient in Equation~\eqref{eq:main-expansion} admits the canonical source-wise decomposition induced by the exact local second-moment identity in Equation~\eqref{eq:D2-exact}:
\begin{align}
  B_{20}^{\mathrm{local}}(H)
  &:=-\frac{\tau\sigma^2}{4a}(H-1),                                  \label{eq:B20local}\\
  B_{20}^{\mathrm{ctrl}}(H)
  &:=-\frac{\tau\sigma^2}{12a}
  \frac{(H-1)(2H-1)}{H},                                               \label{eq:B20ctrl}\\
  B_{20}^{\mathrm{SCAF}}(H)
  &:=B_{20}^{\mathrm{local}}(H)+B_{20}^{\mathrm{ctrl}}(H)
  =-\frac{\tau\sigma^2}{12a}
  \frac{(H-1)(5H-1)}{H}.                                               \label{eq:B20total}
\end{align}
\end{corollary}

\begin{proof}
The direct fresh-noise source in Equation~\eqref{eq:D2-exact} is $\gamma^2h\sigma^2$. Its contribution through the factor $-\tau/(2aH)$ in Equation~\eqref{eq:b-assembly} is
\[
  -\frac\tau{2aH}\sigma^2\sum_{h=0}^{H-1}h
  =-\frac{\tau\sigma^2}{4a}(H-1).
\]
The $N^0$ part of the control source is
\[
  \gamma^2h^2\frac{\sigma^2}{H}.
\]
Therefore, its contribution is
\[
  -\frac\tau{2aH}\frac{\sigma^2}{H}
  \sum_{h=0}^{H-1}h^2
  =-\frac{\tau\sigma^2}{12a}
  \frac{(H-1)(2H-1)}{H}.
\]
Adding the two gives Equation~\eqref{eq:B20total}.
\end{proof}

\begin{remark}[Meaning of the decomposition]
The word \emph{canonical} refers to the two distinct terms in the exact square in Equation~\eqref{eq:D2-exact}: direct within-round fresh noise and the persistent round-start control second moment. We do not claim uniqueness under arbitrary algebraic reparameterizations. We also do not identify $B_{20}^{\mathrm{local}}$ with a complete general-$H$ FedAvg coefficient; that would require a separate FedAvg stationary theorem or comparison proof.
\end{remark}

\begin{remark}[Finite-$N$ meaning of $B_{20}^{\mathrm{ctrl}}$]
Equation~\eqref{eq:B20ctrl} is the client-number-independent $N^0\gamma^2$ coefficient induced by the control second moment. Because Equation~\eqref{eq:QH-uniform} contains a finite-$N$ correction proportional to $1/N$, the complete finite-$N$ control-generated $\gamma^2$ contribution may contain additional $\gamma^2/N$ terms, all retained in the remainder in Equation~\eqref{eq:main-remainder}.
\end{remark}

\section{Uniform constants and joint-limit semantics}
\label{sec:uniformity}

\subsection{Uniformity in \texorpdfstring{$N$ and $\gamma$}{N and gamma} for fixed \texorpdfstring{$H$}{H}}

\begin{proposition}[Common small-step threshold and uniform remainder constant]
\label{prop:uniformity}
For fixed $H$, all lemmas used in \cref{thm:main} admit a common threshold $\gamma_{0, H}>0$ and constants that are independent of $N$ and $\gamma$. Consequently, the remainder in Equation~\eqref{eq:main-remainder} is uniform over
\[
  N\ge2,\qquad 0<\gamma\le\gamma_{0,H}.
\]
\end{proposition}

\begin{proof}
Each preceding argument requires only finitely many small-step restrictions of the form $\gamma\le c_H$, where $c_H>0$ depends on fixed $H$ and fixed problem/noise parameters.
These include the stationarity and moment restrictions inherited from \citet{mangold2025scaffold}, $\gamma\le1$, conditions such as $a\gamma\le1$ or $\gamma HL\le1$, finite-step stability restrictions, and the fixed-$\epsilon$ absorption conditions in \cref{lem:M4}. Define
\[
  \gamma_{0,H}:=\min_j\gamma_{0,H}^{(j)}.
\]
Because the collection is finite and every $\gamma_{0, H}^{(j)}>0$, the common threshold is positive and independent of $N$.

The constant bookkeeping is uniform in $N$ for the following reasons.
\begin{itemize}
  \item The imported coarse iterate moments and the restricted control expansion \cref{lem:QH-uniform} have constants independent of $N$.
  \item Client averaging introduces factors $N^{-1}$ or $N^{-2}$ explicitly; Jensen and Cauchy--Schwarz do not introduce hidden powers of $N$.
  \item Mixed fractional scales are reduced using algebraic inequalities such as
  \[
    \frac{\gamma^{5/2}}{\sqrt N}
    \le \frac12\left(\frac{\gamma^2}{N}+\gamma^3\right),
  \]
  with constants independent of the relative rate of $N$ and $\gamma$.
  \item In the coordinate-replacement argument in Equation~\eqref{eq:etaDelta-sharp}, the $NH$ fresh coordinates are multiplied by one coefficient of order $\gamma/N$ and one path-sensitivity factor of order $\gamma^2/N$, leaving $C_H\gamma^3/N$, not an $N^0$ term.
  \item Restoring factors satisfy lower bounds of order $a\gamma$, and $a\ge\mu>0$ is fixed; the resulting division by $\gamma$ is matched by an extra factor of $\gamma$ in the numerator of each moment estimate.
  \item The fourth-moment absorption fixes $\epsilon>0$ after the $H$-dependent constants are known; $\epsilon$ is never chosen as a function of $N$ or $\gamma$.
\end{itemize}
Thus the final constant $C_H$ may depend on $H,\mu,L,K_3,K_4,K_5,B_\eps,\sigma^2$ and related fixed quantities, but not on $N$ or $\gamma$.
\end{proof}

\begin{remark}[No explicit closed form for $\gamma_{0,H}$ is claimed]
The theorem requires only the existence of a positive common threshold. Writing a closed-form minimum of all numerical restrictions would require a separate mechanical constant chase and is intentionally omitted. This avoids presenting a false ``explicit'' threshold that silently fails to meet an absorption condition or a restriction inherited from the earlier analysis.
\end{remark}

\begin{remark}[Not uniform in $H$]
Neither $C_H$ nor $\gamma_{0,H}$ is claimed to be controlled uniformly as $H\to\infty$. Factors generated by finite-step stability, sums over local steps, and the coefficient itself are allowed to grow with $H$. The theorem is an arbitrary-fixed-$H$ statement only.
\end{remark}

\subsection{Uniform joint \texorpdfstring{$(\gamma,1/N)$}{(gamma, 1/N)} interpretation}

\begin{corollary}[Joint-limit characterization of the $N^0\gamma^2$ coefficient]
\label{cor:joint-limit}
For any sequences $\gamma_k\to0$ and integers $N_k\to\infty$, with fixed $H$ and no relative-rate condition,
\begin{equation}
\label{eq:joint-limit}
  \lim_{k\to\infty}
  \frac{
    b_{\gamma_k,N_k,H}
    +\dfrac{\tau\sigma^2}{4a^2}\dfrac{\gamma_k}{N_k}
  }{\gamma_k^2}
  =-\frac{\tau\sigma^2}{12a}
  \frac{(H-1)(5H-1)}{H}.
\end{equation}
\end{corollary}

\begin{proof}
Divide Equation~\eqref{eq:main-remainder} by $\gamma^2$:
\[
  \left|\frac{R_{\gamma,N,H}}{\gamma^2}\right|
  \le C_H\left(\frac1N+\gamma\right).
\]
Both terms converge to zero along any such sequence, independently of their relative rate.
\end{proof}

\begin{remark}[Correct asymptotic terminology]
The monomials $\gamma/N$ and $\gamma^2$ do not have a fixed scalar ordering along all two-parameter paths. For example, $N=\gamma^{-3}$ makes $\gamma/N=\gamma^4$, while $N=\gamma^{-1/2}$ makes $\gamma/N=\gamma^{3/2}$. Therefore, the precise statement supported by \cref{thm:main} is that \cref{thm:main} identifies the \emph{client-number-independent $N^0\gamma^2$ coefficient in a uniform joint $(\gamma,1/N)$ expansion after removing the known $\gamma/N$ layer}. It should not be described unconditionally as ``the next term.''
\end{remark}

\begin{remark}[When the coefficient is nonzero]
The coefficient in Equation~\eqref{eq:B20total} is nonzero whenever
\[
  \tau\sigma^2(H-1)\ne0.
\]
Without this additional condition, the theorem identifies the $N^0\gamma^2$ coefficient but does not assert that it is the first nonzero client-independent term.
\end{remark}

\subsection{Precise claim boundary}

The strongest claim supported by the present proof is:

\begin{quote}
For fixed $H$, we derive a stationary-bias expansion for one-dimensional homogeneous stochastic SCAFFOLD that is uniform jointly in $(\gamma,1/N)$. After removing the known $O(\gamma/N)$ layer, the client-number-independent $N^0\gamma^2$ coefficient is
\[
  -\frac{f'''(x^\star)\sigma^2}{12f''(x^\star)}
  \frac{(H-1)(5H-1)}{H}.
\]
The coefficient admits a canonical source-wise decomposition into a direct local-gradient-noise contribution and an additional contribution induced by the persistent SCAFFOLD control second moment.
\end{quote}

The following claims are \emph{not} established here:
\begin{itemize}
  \item a complete second-order expansion at fixed finite $N$;
  \item discovery of the fact that stochastic SCAFFOLD is biased;
  \item a claim that all higher-order bias is caused by control variates;
  \item uniformity as $H\to\infty$;
  \item heterogeneous-client, multidimensional, or state-dependent-noise extensions;
  \item equality of $B_{20}^{\mathrm{local}}$ with a separately proved general-$H$ FedAvg coefficient.
\end{itemize}

\section{Numerical protocol and supplementary checks}
\label{sec:numerical-supp}

This appendix records the simulation protocol, the initial coefficient grid, and the negative controls summarized in \cref{sec:numerics}.

\subsection{Full simulation protocol}

For each setting, we estimate the stationary mean in two ways: by the direct sample mean and by the exact mean-balance identity used in the proof. Unless an entry is explicitly labeled ``direct,'' the reported $\hat b_{\gamma, N,H}$ and normalized statistic $C_{\gamma, N,H}$ use the mean-balance estimator; the direct sample mean is retained as a diagnostic. Let
\[
m_{\gamma,H}=\left\lceil\frac{1}{\mu H\gamma}\right\rceil.
\]
We discard $100m_{\gamma, H}$ communication rounds as burn-in and use batches of length $20m_{\gamma, H}$, with at least 100 batches per chain. Each setting uses eight independent chains. Pointwise uncertainty is reported using 95\% Student-$t$ intervals across the eight chain estimates.

For extrapolated intercepts, we use a nonparametric bootstrap at the independent-chain level. Within each parameter setting, the eight chain estimates are resampled with replacement, the setting mean is recomputed, and the same unweighted linear regression is refit. The reported 95\% bootstrap interval is the interval between the empirical 2.5\% and 97.5\% quantiles of the refitted intercepts.

For every reported setting, we monitor finiteness, preservation of the zero-sum control invariant, agreement between the direct and mean-balance estimators, and stability between the first and second halves of the post-burn-in sample. These checks diagnose simulation failures or insufficient equilibration; they are not tests of the theorem.

\subsection{Initial joint path and the recorded \texorpdfstring{$H=4$}{H=4} discrepancy}

The initial study followed $N=1/\gamma$ at
\[
\gamma\in\{1/12,1/16,1/24,1/32,1/48\},
\qquad H\in\{2,4\}.
\]
For $H=2$, a linear extrapolation gives intercept $-0.187214$ with 95\% bootstrap interval $[-0.189250,-0.185175]$, consistent with $B_{20}(2)=-0.1875$. For $H=4$, the same coarse-grid extrapolation gives $-0.588633$ with interval $[-0.590776,-0.586514]$, which does not contain $B_{20}(4)=-0.59375$. Under the criterion specified for the initial study, this outcome was recorded as a contradiction candidate. The theorem permits an $O_H(\gamma)$-normalized remainder and does not imply exact linearity over this finite grid, so we specified a smaller-step experiment with fresh random seeds before inspecting its outcomes.

\begin{figure}[H]
\centering
\begin{subfigure}{0.49\linewidth}
  \centering
  \includegraphics[width=\linewidth]{figures/coefficient_convergence_H2.png}
  \caption{$H=2$.}
\end{subfigure}\hfill
\begin{subfigure}{0.49\linewidth}
  \centering
  \includegraphics[width=\linewidth]{figures/coefficient_convergence_H4.png}
  \caption{$H=4$.}
\end{subfigure}
\caption{Initial coarse-grid coefficient study. The $H=2$ extrapolation is consistent with the target coefficient, whereas the $H=4$ extrapolation is offset on this finite-step grid.}
\label{fig:initial-coeff}
\end{figure}

The follow-up used
\[
(\gamma,N)\in\{(1/64,64),(1/96,96),(1/128,128),(1/192,192)\}.
\]
Writing $E(\gamma):=C_{\gamma,1/\gamma,4}-B_{20}(4)$, the observed errors decrease from $0.02159$ to $0.00787$, while $E(\gamma)/\gamma$ remains of constant order. The follow-up met its specified validity, directional convergence, smallest-step proximity, and $O(\gamma)$-compatibility criteria.

\begin{figure}[H]
\centering
\begin{subfigure}{0.49\linewidth}
  \centering
  \includegraphics[width=\linewidth]{figures/coefficient_vs_step_size.png}
  \caption{Normalized residual.}
\end{subfigure}\hfill
\begin{subfigure}{0.49\linewidth}
  \centering
  \includegraphics[width=\linewidth]{figures/scaled_error_vs_step_size.png}
  \caption{Error divided by $\gamma$.}
\end{subfigure}
\caption{Smaller-step $H=4$ study. The normalized residual moves toward the predicted coefficient, while the scaled error remains of constant order, consistent with the \(O_H(\gamma)\) remainder permitted by Equation~\eqref{eq:Csim-prediction} along \(N=1/\gamma\).}
\label{fig:H4-confirm}
\end{figure}

\begin{table}[H]
\centering
\caption{Fresh smaller-step results for $H=4$. Intervals are 95\% Student-$t$ intervals across eight independent chains.}
\label{tab:H4-small}
\begin{tabular}{@{}rrrrr@{}}
\toprule
$\gamma$ & $N$ & $C_{\gamma,N,4}$ & 95\% CI & $E(\gamma)/\gamma$\\
\midrule
$1/64$  & 64  & $-0.572157$ & $[-0.575671,-0.568643]$ & $1.382$\\
$1/96$  & 96  & $-0.579148$ & $[-0.581447,-0.576848]$ & $1.402$\\
$1/128$ & 128 & $-0.582795$ & $[-0.584912,-0.580678]$ & $1.402$\\
$1/192$ & 192 & $-0.585876$ & $[-0.587794,-0.583958]$ & $1.512$\\
\bottomrule
\end{tabular}
\end{table}

\subsection{Negative controls and the \texorpdfstring{$H=1$}{H=1} boundary}

With deterministic gradients ($\sigma^2=0$), both displayed stochastic coefficients vanish. For the quadratic objective $f'(x)=x$, one has $f'''(x^\star)=0$; the mean-balance estimator is then structurally zero, so we report the direct sample mean. Finally, $H=1$ is a boundary check outside the theorem's stated $H\ge2$ domain. The statistic in that case is the normalized residual after removing the leading $\gamma/N$ term.

\begin{table}[H]
\centering
\small
\caption{Negative controls and the $H=1$ boundary check.}
\label{tab:negative-controls}
\begin{tabular}{@{}lrrcc@{}}
\toprule
Case & $\gamma$ & $N$ & Statistic & 95\% CI\\
\midrule
Deterministic, $H=2$ & $1/24$ & 24 & $\hat b_{\rm direct}=5.39\times10^{-16}$ & $[-1.49\times10^{-16},\,1.23\times10^{-15}]$\\
Deterministic, $H=2$ & $1/48$ & 48 & $\hat b_{\rm direct}=-3.57\times10^{-15}$ & $[-1.06\times10^{-14},\,3.47\times10^{-15}]$\\
Quadratic, $H=2$ & $1/24$ & 24 & $\hat b_{\rm direct}=7.11\times10^{-4}$ & $[-1.98\times10^{-4},\,1.62\times10^{-3}]$\\
Quadratic, $H=2$ & $1/48$ & 48 & $\hat b_{\rm direct}=1.92\times10^{-4}$ & $[-3.03\times10^{-4},\,6.86\times10^{-4}]$\\
$H=1$ boundary & $1/24$ & 24 & $C=-2.345\times10^{-3}$ & $[-4.156\times10^{-3},\,-5.33\times10^{-4}]$\\
$H=1$ boundary & $1/48$ & 48 & $C=-4.87\times10^{-4}$ & $[-3.658\times10^{-3},\,2.684\times10^{-3}]$\\
\bottomrule
\end{tabular}
\end{table}

The deterministic and quadratic intervals contain zero at both tested step sizes. For $H=1$, the interval is slightly separated from zero at $\gamma=1/24$ but contains zero at $\gamma=1/48$, with a smaller point estimate at the smaller step size.
\end{document}
